% \iffalse meta-comment
% !TEX program  = pdfLaTeX
% %%%%%%%%%%%%%%%%%%%%%%%%%%%%%%%%%%%%%%%%%%%%%%%%%%%%%%%%%%%%%%%%%%%%%%%%%%%
%                                                                           %
%    CHANGE THESE VALUES WITH EACH NEW RELEASE:                             %
%                                                                           %
%<class>\ProvidesClass{apa7}[2022/12/03 v2.17 APA formatting (7th edition)]
%                                                                           %
%<*internal>                                                                %
\def\apaSevenVersionDate{2022/12/03}
\def\apaSevenVersionNumber{2.17}
%                                                                           %
%                                                                           %
% %%%%%%%%%%%%%%%%%%%%%%%%%%%%%%%%%%%%%%%%%%%%%%%%%%%%%%%%%%%%%%%%%%%%%%%%%%%
\iffalse
%</internal>
%<*readme>
----------------------------------------------------------------------

apa7 - A LaTeX class for formatting documents in compliance with the
American Psychological Association's Publication Manual, 7th edition

Copyright (C) 2022 by Daniel A. Weiss <daniel.weiss.led at gmail.com>

This work may be distributed and/or modified under the
conditions of the LaTeX Project Public License (LPPL), either
version 1.3c of this license or (at your option) any later
version.  The latest version of this license is in the file:

http://www.latex-project.org/lppl.txt

Users may freely modify these files without permission, as long as the
copyright line and this statement are maintained intact.

This work is not endorsed by, affiliated with, or probably even known
by, the American Psychological Association.

----------------------------------------------------------------------

The Publication Manual of the American Psychological Association is
widely used in the social sciences.  The most recent update, in 2019,
altered the formatting guidelines and therefore rendered existing
formatting solutions inadequate.  The apa7 class is an update of
older code from the apa6 class (available from CTAN).  The student 
manuscript type has been added. See the documentation for details.

----------------------------------------------------------------------

Installation Notes:
1. Create directory "config" next to the file apa7.ins
2. Run latex apa7.ins
3. Move/copy apa7.cls to the directory of your latex install or the directory containing your latex project.
%</readme>
%<*internal>
\fi
\def\nameofplainTeX{plain}
\ifx\fmtname\nameofplainTeX\else
  \expandafter\begingroup
\fi
%</internal>
%<*install>
\input docstrip.tex
\keepsilent
\askforoverwritefalse
\preamble
----------------------------------------------------------------------

apa7 - A LaTeX class for formatting documents in compliance with the
American Psychological Association's Publication Manual, 7th edition

Copyright (C) 2022 by Daniel A. Weiss <daniel.weiss.led at gmail.com>

This work may be distributed and/or modified under the
conditions of the LaTeX Project Public License (LPPL), either
version 1.3c of this license or (at your option) any later
version.  The latest version of this license is in the file:

http://www.latex-project.org/lppl.txt

Users may freely modify these files without permission, as long as the
copyright line and this statement are maintained intact.

This work is not endorsed by, affiliated with, or probably even known
by, the American Psychological Association.

----------------------------------------------------------------------

\endpreamble
\postamble

Copyright (C) 2022 by Daniel A. Weiss <daniel.weiss.led at gmail.com>

This work may be distributed and/or modified under the
conditions of the LaTeX Project Public License (LPPL), either
version 1.3c of this license or (at your option) any later
version.  The latest version of this license is in the file:

http://www.latex-project.org/lppl.txt

Users may freely modify these files without permission, as long as the
copyright line and this statement are maintained intact.

This work is not endorsed by, affiliated with, or probably even known
by, the American Psychological Association.


This work is "maintained" (as per LPPL maintenance status) by
Daniel A. Weiss.

This work consists of the file  apa7.dtx
and the derived files           apa7.ins,
                                apa7.cls,
                                apa7.pdf,
                                README,
                                apa7-american.trsl,
                                apa7-british.trsl,
                                apa7-czech.trsl,
                                apa7-dutch.trsl,
                                apa7-english.trsl,
                                apa7-french.trsl,
                                apa7-german.trsl,
                                apa7-greek.trsl,
                                apa7-hungarian.trsl,
                                apa7-ngerman.trsl,
                                apa7-spanish.trsl,
                                apa7-turkish.trsl,
                                APA7endfloat.cfg,
                                Figure1.pdf,
                                shortsample.tex,
                                longsample.tex, and
                                bibliography.bib.

\endpostamble
\usedir{tex/latex/apa7}
\generate{
  \file{\jobname.cls}{\from{\jobname.dtx}{class}}
  \file{apa7-american.trsl}{\from{\jobname.dtx}{american}}
  \file{apa7-british.trsl}{\from{\jobname.dtx}{british}}
  \file{apa7-czech.trsl}{\from{\jobname.dtx}{czech}}
  \file{apa7-dutch.trsl}{\from{\jobname.dtx}{dutch}}
  \file{apa7-english.trsl}{\from{\jobname.dtx}{english}}
  \file{apa7-french.trsl}{\from{\jobname.dtx}{french}}
  \file{apa7-german.trsl}{\from{\jobname.dtx}{german}}
  \file{apa7-greek.trsl}{\from{\jobname.dtx}{greek}}
  \file{apa7-hungarian.trsl}{\from{\jobname.dtx}{hungarian}}
  \file{apa7-ngerman.trsl}{\from{\jobname.dtx}{ngerman}}
  \file{apa7-spanish.trsl}{\from{\jobname.dtx}{spanish}}
  \file{apa7-turkish.trsl}{\from{\jobname.dtx}{turkish}}
}
\usedir{tex/latex/apa7/config}
\generate{
  \file{./config/APA7endfloat.cfg}{\from{\jobname.dtx}{APAendfloat}}
}
\usedir{tex/latex/apa7/samples}
\generate{
  \file{./samples/bibliography.bib}{\from{\jobname.dtx}{bibliography}}
  \file{./samples/shortsample.tex}{\from{\jobname.dtx}{shortsample}}
  \file{./samples/longsample.tex}{\from{\jobname.dtx}{longsample}}
}
%</install>
%<install>\endbatchfile
%<*internal>
\usedir{source/latex/apa7}
\generate{
  \file{\jobname.ins}{\from{\jobname.dtx}{install}}
}
\nopreamble\nopostamble
\usedir{doc/latex/apa7}
\generate{
  \file{README.txt}{\from{\jobname.dtx}{readme}}
}
\ifx\fmtname\nameofplainTeX
  \expandafter\endbatchfile
\else
  \expandafter\endgroup
\fi
%</internal>
%<class>\NeedsTeXFormat{LaTeX2e}
%<*driver>
\documentclass{ltxdoc}
\usepackage[T1]{fontenc}
\usepackage{lmodern}
\RequirePackage{scalerel,tikz} % this is included for ORCID icon
\usetikzlibrary{svg.path}
\definecolor{orcidlogocol}{HTML}{A6CE39}
\tikzset{
  orcidlogo/.pic={
    \fill[orcidlogocol] svg{M256,128c0,70.7-57.3,128-128,128C57.3,256,0,198.7,0,128C0,57.3,57.3,0,128,0C198.7,0,256,57.3,256,128z};
    \fill[white] svg{M86.3,186.2H70.9V79.1h15.4v48.4V186.2z}
                 svg{M108.9,79.1h41.6c39.6,0,57,28.3,57,53.6c0,27.5-21.5,53.6-56.8,53.6h-41.8V79.1z M124.3,172.4h24.5c34.9,0,42.9-26.5,42.9-39.7c0-21.5-13.7-39.7-43.7-39.7h-23.7V172.4z}
                 svg{M88.7,56.8c0,5.5-4.5,10.1-10.1,10.1c-5.6,0-10.1-4.6-10.1-10.1c0-5.6,4.5-10.1,10.1-10.1C84.2,46.7,88.7,51.3,88.7,56.8z};
  }
}
\newcommand{\addORCIDlink}[2]{#1 \href{https://orcid.org/#2}{{\mbox{\scalerel*{
\begin{tikzpicture}[yscale=-1,transform shape]
\pic{orcidlogo};
\end{tikzpicture}
}{|}}} https://orcid.org/#2}}

\usepackage{hyperref}
\hypersetup{%
   colorlinks,
   linkcolor=blue,
   citecolor=blue,
   urlcolor=black,
   pdfauthor   = {Daniel A. Weiss <daniel.weiss.led@gmail.com>},
   pdftitle    = {The apa7 class},
   pdfsubject  = {Documentation of LaTeX class apa7},
   pdfkeywords = {LaTeX, APA}
}
\usepackage{graphicx}
\usepackage{booktabs}
\usepackage{verbatimbox}
\EnableCrossrefs
\CodelineIndex
\OnlyDescription  % don't publish the full code in the PDF documentation file
\RecordChanges
\begin{document}
  \DocInput{\jobname.dtx}
\end{document}
%</driver>
% \fi
%
%\DoNotIndex{\#,\$,\%,\&,\@,\\,\{,\},\^,\_,\~,\ }
%\DoNotIndex{\@ne}
%\DoNotIndex{\advance,\begingroup,\catcode,\closein}
%\DoNotIndex{\closeout,\day,\def,\edef,\else,\empty,\endgroup}
%\DoNotIndex{\newcommand,\renewcommand}
%
%
% \title{^^A
%     Formatting documents in APA style (7th Edition)\\^^A
%     with the \textsf{apa7} \LaTeX\ class^^A
%     \thanks{This file describes version ^^A
%     \apaSevenVersionNumber, last revised \apaSevenVersionDate.}^^A
% }
% \author{Daniel A. Weiss\thanks{E-mail: daniel.weiss.led@gmail.com}}
% \date{Released \apaSevenVersionDate}
%
% \maketitle
%
% \changes{v1.00}{2019/12/26}{Initial release}
%
% \changes{v1.01}{2019/12/29}{Corrected figure caption styles for doc format.}
% \changes{v1.01}{2019/12/29}{Clarified that apacite and natbib are in 6th edition style.}
%
% \changes{v1.02}{2019/12/29}{Fixed filename issue with language localization files.}
%
% \changes{v1.03}{2019/12/29}{Updated class to detect babel package in document and localize accordingly.}
%
% \changes{v1.04}{2019/12/30}{Added table of contents. Updated documentation and added correct license statement to cls for apa6 package.}
%
% \changes{v1.05}{2020/03/27}{Fixed float placement at end for man and stu mode to ensure one correctly placed figure/table per page. Thanks WC.}
%
% \changes{v1.06}{2020/05/04}{Added figurenote command.}
%
% \changes{v1.07}{2020/05/10}{Fixed bug in figurenote command. Fixed figure title formatting in doc mode. Thanks RH.}
%
% \changes{v1.08}{2020/05/27}{Added support for optional running head (shorttitle) in stu mode.}
%
% \changes{v1.09}{2020/06/19}{Fixed note command (will no longer print "Author Note" header).}
%
% \changes{v2.01}{2020/07/18}{Re-wrote author and affiliation to be APA7 compliant. Fixed spacing on Title Page for title and author. }
%
% \changes{v2.02}{2020/08/02}{Added the "and" in multiple author list to localization files. }
%
% \changes{v2.03}{2020/08/12}{Added French localization. }
%
% \changes{v2.04}{2020/10/11}{Updated ORCID link style for jou mode. }
%
% \changes{v2.05}{2020/11/14}{Fixed a bug with the noextraspace option and added donotrepeattitle option to the documentation. }
%
% \changes{v2.06}{2020/12/02}{Added tablenote command. Added "Note." for table/figure notes to localization files. }
%
% \changes{v2.07}{2020/12/05}{Extended mask command to also mask in stu mode. }
%
% \changes{v2.08}{2021/03/01}{Updated BibLaTex to allow overriding bibliography title. }
%
% \changes{v2.09}{2021/03/03}{Updated APA7endfloat.cfg file to use more modern DeclareDelayedFloatFlavor. This fixes the bug where longtable can be shown on the same page as another table. }
%
% \changes{v2.10}{2021/03/04}{Update BibLaTex instructions to remove DeclareLanguageMapping and remove commands that are no longer required. }
%
% \changes{v2.11}{2021/04/05}{Updated linebreaks in caption separators to prevent warning. }
%
% \changes{v2.12}{2021/04/08}{Reset equation counter for every appendix (like for figures and tables). Thanks mmrabe. }
%
% \changes{v2.13}{2021/05/17}{Fix headers for jou style when journal command is provided. }
%
% \changes{v2.14}{2021/08/08}{Added Spanish localization. Fixed compatibility with Version 4.0 of flushend.sty. Clarify installation in readme. }
%
% \changes{v2.15}{2022/01/03}{Spacing updated to more closely match APA 7th Edition Manual.}
%
% \changes{v2.16}{2022/07/25}{Add Hungarian localization.}
%
% \changes{v2.17}{2022/12/03}{Use tracklang and translations packages to provide translations. Also update Dutch translation.}
%
% \begin{abstract}
%   The \textit{Publication Manual} of the American Psychological
%   Association is widely used in the social sciences.  The most
%   recent update, in 2019, altered the formatting guidelines and
%   therefore rendered existing formatting solutions inadequate.  The
%   \textsf{apa7} class is an update of older code from the
%   \textsf{apa6} class.  The student manuscript type has been added.
% \end{abstract}
%
% \tableofcontents
%
% \section{Background}
% Most journals in the social sciences require manuscripts to be
% formatted in compliance with the American Psychological
% Association's \textit{Publication Manual}, which is updated
% periodically.  The 7th Edition, released in 2019, substantially
% changed the guidelines for formatting manuscripts; these
% modifications rendered existing formatting solutions (e.g., the
% \textsf{apa6} \LaTeX\ class) inadequate for venues in which 7th
% Edition guidelines are being enforced.  The \textsf{apa7} class
% solves this problem, and provides some new functionality not offered
% by the \textsf{apa6} class.
%
% \section{Disclaimer}
% Great care has been taken to ensure the closest possible match
% between APA requirements and the output of this class.  However, it
% is the sole responsibility of the user to ensure compliance with
% specific journal submission requirements!
%
% \section{\textsf{apa7} Version 2.01 Changes to Authors and Affiliations}
% The way that authors and affiliations work has been changed in 
% version 2.01 to be APA 7th edition compliant. 
% This change will require updating any multiple author and/or 
% multiple affiliation documents to use the new |\authorsnames| and
% |\authorsaffiliations| commands (See Section~\ref{sec:authorsaffiliations}).
%
% \section{Usage}
% \label{sec:usage}
%
% \subsection{Class Options}
% \label{sec:classoptions}
% When loading \textsf{apa7} with
% |\documentclass|\oarg{options}|{apa7}|, the following options are
% available.
%
% \vspace{1em}
% \textbf{Document mode:} Four choices are available.
% \begin{itemize}
% \item \DescribeMacro{jou}|jou| (default): Formats the document with an
%   appearance resembling a printed APA journal (e.g., \textit{Journal
%   of Educational Psychology}.  The text is typeset in two-sided,
%   two-column format.
% \item \DescribeMacro{man}|man|: Formats the document in close (if
%   not complete) compliance with the requirements for submission to
%   an APA journal (e.g., title page, double-spacing, etc.).
% \item \DescribeMacro{stu}|stu|: Formats the document in close (if
%   not complete) compliance with the requirements for student paper
%  (e.g., title page, double-spacing, etc.).
% \item \DescribeMacro{doc}|doc|: Formats the document as a typical
%   \LaTeX\ document (one-sided, single-column, etc.)
% \end{itemize}
%
% \textbf{Other class options:}
% \begin{itemize}
% \item \DescribeMacro{10pt}|10pt|: Typesets the document in 10-point
%   font.
% \item \DescribeMacro{11pt}|11pt|: Typesets the document in 11-point
%   font.
% \item \DescribeMacro{12pt}|12pt|: Typesets the document in 12-point
%   font.
% \item \DescribeMacro{a4paper}|a4paper|: Specifies A4 paper size
%   (letter is default).
% \item \DescribeMacro{nolmodern}|nolmodern|: Suppresses loading of the
%   \textsf{lmodern} package.
% \item \DescribeMacro{nofontenc}|nofontenc|: Suppresses loading of the
%   \textsf{fontenc} package (which is needed for proper hyphenation
%   of accented characters).
% \item \DescribeMacro{babel}|babel|: In all modes, loads
%   \textsf{babel}; the desired language(s) are listed as options
%   immediately following |babel|; the last language listed is the
%   main one.
% \item \DescribeMacro{donotrepeattitle}|donotrepeattitle|: Suppress the title above the 
%   introduction (in case a different title is desired).
% \item \DescribeMacro{noextraspace}|noextraspace|: Removed option in version 2.15
%  in an attempt to more closely match APA 7th edition guidelines (will throw error if used).
% \item \DescribeMacro{floatsintext}|floatsintext|: In |man| and |stu| mode,
%   integrates floats (tables and figures) within the body of the text
%   instead of postponing them until after the reference list.
% \item \DescribeMacro{biblatex}|biblatex|: Loads the
%   \textsf{biblatex} package; see Section~\ref{sec:biblatexinfo} for
%   details.
% \item \DescribeMacro{apacite}|apacite|: Loads the \textsf{apacite}
%   package; see Section~\ref{sec:apaciteinfo} for details. \\
%   \textit{\textbf{Note:} at this time, this will produce citations in 6th edition format. }
% \item \DescribeMacro{natbib}|natbib|: See
%   Section~\ref{sec:natbibinfo} for details. \\
%   \textit{\textbf{Note:} at this time, this will produce citations in 6th edition format.}
% \item \DescribeMacro{mask}|mask|: Masks references that are marked
%   as the author's own (for masked peer review); see
%   Section~\ref{sec:masking} for details.
% \item \DescribeMacro{longtable}|longtable|: If you \textit{must} use
%   long tables (exceeding one page in length), try this option (but
%   it may not work in all contexts).  Do \textbf{\textit{not}} load
%   \textsf{longtable} yourself because of precedence requirements
%   with the \textsf{endfloat} package.  Copy the file
%   |APA7endfloat.cfg| from the ``config'' folder of your \textsf{apa}
%   installation to the working folder of your document (not in your
%   |texmf| tree), and rename it to |endfloat.cfg| so that endfloat
%   will recognize it.  The supplied |APA7endfloat.cfg| file will also
%   be necessary in conjunction with the \textsf{rotating} package
%   (and its |\sideways| command) to produce rotated tables; this
%   works for |man| and |stu| mode only.  For rotated tables in |jou| or |doc|
%   mode, the \textsf{rotating} package may be used.  If sideways
%   tables get pushed to the end of your document with the
%   |floatsintext| option, try using the |ph!| placement specifier for
%   your sideways table(s); for example, |\begin{sidewaystable}[ph!]|.
% \item \DescribeMacro{notxfonts}|notxfonts|: In |jou| mode, prevents
%   \textsf{txfonts} from loading, in case \textsf{pslatex} or
%   \textsf{times} is preferable for some reason.
% \item \DescribeMacro{notimes}|notimes|: In |jou| mode, cancels
%   loading \textsf{txfonts} or \textsf{pslatex} or \textsf{times} and
%   uses Computer Modern instead.
% \item \DescribeMacro{notab}|notab|: In |jou| mode, cancels the
%   automatic stretching of tabular environments to the width of their
%   enclosing float.
% \item \DescribeMacro{helv}|helv|: In |man| and |stu| mode,
%   uses Helvetica font instead of Computer Modern.
% \item \DescribeMacro{nosf}|nosf|: In |man| and |stu| mode,
%   neutralizes the |\helvetica| command.
% \item \DescribeMacro{tt}|tt|: In |man| and |stu| mode, uses typewriter-like
%   font.
% \item \DescribeMacro{draftfirst}|draftfirst|: In all modes, places
%   the word ``DRAFT'' as a watermark across the first page.
% \item \DescribeMacro{draftall}|draftall|: In all modes, places the
%   word ``DRAFT'' as a watermark across all pages.
% \end{itemize}
%
% Class options not handled by \textsf{apa7} (e.g., |draft|) will be
% passed on to the \textsf{article} class.
%
% \subsection{Document Preamble}
% \label{sec:preamble}
% The following commands are available within the document preamble
% (i.e., the part of the file preceding |\begin{document}|).
%
% \subsubsection{Document Title and Headers}
% \label{sec:titleheaders}
% \begin{itemize}
% \item \DescribeMacro{\title}|\title|\marg{document-title}: The title of the document
% \item \DescribeMacro{\shorttitle}|\shorttitle|\marg{short-title}: A
%   shortened version of the title (for page headers)
% \item
%   \DescribeMacro{\leftheader}|\leftheader|\marg{author-last-name(s)}:
%   Author last name(s) (for even-page headers in |jou| mode)
% \end{itemize}
% Optional; use if desired:\par
% \begin{itemize}
% \item \DescribeMacro{\journal}|\journal|\marg{journal-name}: Journal
%   name or other note; typeset in the top left header of page 1
%   (|jou| and |doc| modes only); to change the starting page to a
%   number other than 1, insert the following line immediately after
%   |\maketitle|:\\
%   |\setcounter{page}|\marg{custom-page-number}
% \item \DescribeMacro{\volume}|\volume|\marg{journal-volume}: Volume,
%   number, pages; typeset in the top left header in |jou| and |doc|
%   modes, underneath the content of |\journal|
% \item \DescribeMacro{\ccoppy}|\ccoppy|\marg{copright-notice}:
%   Copyright notice, etc.; typeset in the top right header of page 1
%   (|jou| and |doc| modes only)
% \item \DescribeMacro{\copnum}|\copnum|\marg{more-copyright-info}:
%   Any additional text needed; typeset in the top right header in
%   |jou| and |doc| modes, underneath the content of |\ccoppy|
% \end{itemize}
%
% \subsubsection{Authors and Affiliations}
% \label{sec:authorsaffiliations}
% Based on the new APA 7th edition style guide, the previous method for 
% handling multiple authors and multiple affiliations has been deprecated as of version 2.01 of apa7. \par
% \textit{The followings commands have been deprecated and will return an error:} \par
% |\twoauthors{First Author(s)}{Second Author(s)}|\par
% |\threeauthors{John and Jim}{Mary and Sue}{Nick}|\par
% |\fourauthors{Helen}{Dick}{Tracy and Larry}{James Bond}|\par
% |\fiveauthors{...}{...}{...}{...}{...}|\par
% |\sixauthors{...}{...}{...}{...}{...}{...}|\par
% |\twoaffiliations{Affil. of 1st Author(s)}{Affil. 2nd Author(s)}|\par
% |\threeaffiliations{U of A}{U of B}{U of C}|\par
% |\fouraffiliations{My Company}{Your Department}{Heaven}{Earth}|\par
% |\fiveaffiliations{...}{...}{...}{...}{...}|\par
% |\sixaffiliations{...}{...}{...}{...}{...}{...}|\par
% \begin{itemize}
% \item \DescribeMacro{\author}|\author|\marg{author name}: Author name (if only 1 author)
% \item \DescribeMacro{\affiliation}|\affiliation|\marg{author's affiliation}: Author's affiliation (if only 1 affiliation)
% \item \DescribeMacro{\authorsnames}|\authorsnames|\oarg{comma separated list of superscripts}\marg{comma separated list of authors name(s)}: Authors names (if more than 1 author)
% \end{itemize}
%
% Where \marg{authors names} is a comma separated list of authors in display order
% and \oarg{superscripts} is an optional comma separated list of affiliation superscripts. 
% If there are multiple superscripts for one author, put the set of superscripts within |{ }|.
%
% \begin{itemize}
% \item
%   \DescribeMacro{\authorsaffiliations}|\authorsaffiliations|\marg{comma separated list of affiliation(s)}:
%   Author affiliation(s). Put each affiliation within |{ }| and separate each affiliation with a comma.
% \end{itemize}
%
% \noindent \textbf{Examples:} \par
% \noindent One author, one affiliation \par
% |\authorsnames{Maggie C. Leonard}| \par
% |\authorsaffiliations{{Department of Psychology, George Mason University}}| \par
% \noindent One author, two affiliations \par
% |\authorsnames{Andrew K. Jones-Willoughby}| \par
% |\authorsaffiliations{{School of Psychology, University of Sydney},| \par
% |{Center for Behavioral Neuroscience, American University}}| \par
% \noindent Two authors, shared affiliation \par
% |\authorsnames{Mackenzie J. Clement,Talia R. Cummings}| \par
% |\authorsaffiliations{{College of Nursing, Michigan State University}}| \par
% \noindent Two authors, different affiliations \par
% |\authorsnames[1,2]{Wilhelm T. Webber,Latasha P. Jackson}| \par
% |\authorsaffiliations{{Max Planck Institute for Human Development, Berlin, Germany},| \par
% |{College of Education, University of Georgia}}| \par
% \noindent Three or more authors, different affiliations \par
% |\authorsnames[1,{2,3},1]{Savannah C. St. John, Fen-Lei Chang, Carlos O. Vásquez III}| \par
% |\authorsaffiliations{{Education Testing Service, Princeton, New Jersey, ...},| \par
% |{MRC Cognition and Brain Science Unit, Cambridge, England},| \par
% |{Department of Psychology, University of Cambridge}}| \par
%
%
% \subsubsection{Title Page Information}
% \label{sec:titlepageinfo}
% \begin{itemize}
% \item \DescribeMacro{\course}|\course|\marg{course-number-name}: The
%   course number and name (e.g. PSYC 3170: Clinical Psychology)
% \item \DescribeMacro{\professor}|\professor|\marg{professor}: The
%   course professor name.
% \item \DescribeMacro{\duedate}|\duedate|\marg{paper-duedate}: The
%   due date for the paper (e.g. October 1, 2019)
% \item \DescribeMacro{\abstract}|\abstract|\marg{abstract-text}: The
%   abstract of the article
% \item \DescribeMacro{\keywords}|\keywords|\marg{keywords}: Keywords
%   (typeset after the abstract) See Section~\ref{sec:keywords} for
%   details regarding language localization of the ``Keywords'' label.
% \item \DescribeMacro{\authornote}|\authornote|\marg{author-note}:
%   The Author Note, containing contact information, acknowledgements,
%   etc.
% \item \DescribeMacro{\addORCIDlink}|\addORCIDlink|\marg{author-name}\marg{ORCID}: Optionally, use this within the author note to generate the ORCID link. \\
% For example
% \begin{verbatim} \addORCIDlink{First M. Last}{0000-0000-0000-0000} \end{verbatim} \\
% generates \\ \\
% \addORCIDlink{First M. Last}{0000-0000-0000-0000}
% \end{itemize}
%
% Optional; use if desired:\par
% \begin{itemize}
% \item \DescribeMacro{\note}|\note|\marg{note-text}: Notation of
%   manuscript date or other information desired beneath the
%   affiliation line (not part of the APA 7th edition specification for title page)
% \end{itemize}
%
% \subsection{Maketitle}
% The \DescribeMacro{\maketitle}|\maketitle| command formats the
% document title, page headers, author list, author affiliations,
% Author Note (if provided), abstract according to whether |jou|,
% |man|, |stu|, or |doc| mode has been specified.  This command should be
% on the line after |\begin{document}|, with the first line of text
% immediately following the |\maketitle| line (no blank lines).
%
% \subsection{Heading Levels}
% \label{sec:headinglevels}
% Heading levels are automatically formatted using the following
% standard \LaTeX\ commands:
% \begin{itemize}
% \item \DescribeMacro{\section}|\section|\marg{title}
% \item \DescribeMacro{\subsection}|\subsection|\marg{title}
% \item \DescribeMacro{\subsubsection}|\subsubsection|\marg{title}
% \item \DescribeMacro{\paragraph}|\paragraph|\marg{title}
% \item \DescribeMacro{\subparagraph}|\subparagraph|\marg{title}
% \end{itemize}
%
% Please note that sections cannot be |\ref|'d since APA style does
% not use numbered sections.  So |\label| commands are unnecessary
% unless you wish to use |\refname|.
%
% \subsection{Enumeration}
% \label{sec:enumeration}
% Several forms of enumeration are provided, as follows.
%
% \vspace{0.6em}
% \DescribeMacro{\begin\{seriate\}} \DescribeMacro{\end\{seriate\}}
%   |Blah blah blah|\par
%   |\begin{seriate}|\par
%   |  \item first item,|\par
%   |  \item second item.|\par
%   |\end{seriate}|\par
%   |Blah blah blah|\par\vspace{0.6em}
% \noindent results in:\par\vspace{0.6em}
%   |Blah blah blah (a) first item, (b) second item.  Blah blah blah|
%
% \vspace{3em}
% \DescribeMacro{\begin\{APAenumerate\}} \DescribeMacro{\end\{APAenumerate\}}
%   |Blah blah blah|\par
%   |\begin{APAenumerate}|\par
%   |  \item first item ... ... ... continue continue|\par
%   |  \item second item ... ... ... continue continue|\par
%   |\end{APAenumerate}|\par
%   |Blah blah blah|\par\vspace{0.6em}
% \noindent results in:\par\vspace{0.6em}
%   |Blah blah blah|\par
%   |    1. first item ... ... ...|\par
%   |continue continue|\par
%   |    2. second item ... ... ...|\par
%   |continue continue|\par
%   |Blah blah blah|\par
%
% \vspace{3em}
% \DescribeMacro{\begin\{APAitemize\}} \DescribeMacro{\end\{APAitemize\}}
%   |Blah blah blah|\par
%   |\begin{APAitemize}|\par
%   |  \item first item ... ... ... continue continue|\par
%   |  \item second item ... ... ... continue continue|\par
%   |\end{APAitemize}|\par
%   |Blah blah blah|\par\vspace{0.6em}
% \noindent results in:\par\vspace{0.6em}
%   |Blah blah blah|\par
%   |    o first item ... ... ...|\par
%   |continue continue|\par
%   |    o second item ... ... ...|\par
%   |continue continue|\par
%   |Blah blah blah|\par
%
% \vspace{0.6em}
% In addition to the above, all standard \LaTeX\ enumeration
% environments are available (e.g., |enumerate| and |itemize|).
%
% \subsection{Other Macros}
% \label{sec:othermacros}
% \begin{itemize}
% \item \DescribeMacro{\begin\{figure*\}} \DescribeMacro{\end\{figure*\}}
%   \DescribeMacro{\begin\{table*\}} \DescribeMacro{\end\{table*\}}
%   When a figure is too wide for a single column (in |jou| mode), use
%   |\begin{figure*}| and |\end{figure*}| instead of the non-starred
%   version.  The same applies with |\begin{table*}| and |\end{table*}|.
%   When using double-column tables or figures (|jou| mode), use the
%   |\centering| command; for example:\par
%   |\begin{table*}|\par\vspace{-.5em}
%   |  \centering|\par\vspace{-.5em}
%   |    \begin{threeparttable}|\par
% \item
%   \DescribeMacro{\figurenote}|\figurenote|\marg{figure-note-text}:
%   Formats figure note when placed after |\includegraphic|, |\fitfigure|, or similar command within |\begin{figure}| and |\end{figure}|.
% \item
%   \DescribeMacro{\tablenote}|\tablenote|\marg{table-note-text}:
%   Formats table note when placed after |\end{tabular}| or similar command within |\begin{table}| and |\end{table}|.
% \item
%   \DescribeMacro{\fitfigure}|\fitfigure|\oarg{height}\marg{eps-filename}:
%   Automatically fit a postscript figure; use instead of
%   |\includegraphics|
% \item
%   \DescribeMacro{\fitbitmap}|\fitbitmap|\oarg{height}\marg{eps-filename}:
%   Same as |\fitfigure| but won't scale figure in |\man| mode for
%   best reproduction of bitmap figures
% \item \DescribeMacro{\tabfnm}|\tabfnm{a}|: Place a superscript
%   \textbf{f}oot\textbf{n}ote \textbf{m}ark inside a table cell.  Any
%   series of unique identifiers can be used in place of |a|.
% \item \DescribeMacro{\tabfnt}|\tabfnt{a}|\marg{footnote-text}:
%   Within table footnotes, specify the \textbf{f}oot\textbf{n}ote
%   \textbf{t}ext for |\tabfnm{a}|
% \item \DescribeMacro{\apavector}|\apavector|\marg{symbol}: Format
%   the \marg{symbol} as a vector by APA rules
% \end{itemize}
%
% \subsection{Appendices}
% \label{sec:appendices}
% \begin{itemize}
% \item \DescribeMacro{\appendix}|\appendix|: Begins the appendices
%   portion of the document
% \item \DescribeMacro{\section}|\section|\marg{appendix-title}:
%   Begins each appendix
% \end{itemize}
%
% Because appendices are numbered (with letters!) you may establish a
% label for each appendix (e.g., |\label{app:xxx}|); when there is
% more than one appendix, use |Appendix~\ref{app:xxx}| within the main
% body of the text to refer to that appendix.  (Of course, if there is
% only one appendix, simply refer to it as |the Appendix|.)
%
% \section{Known Limitations}
% \begin{itemize}
% \item The |APAenumerate| environment does not nest properly.
% \end{itemize}
%
% \section{Development of \textsf{apa7}}
% The base code for this class is the \textsf{apa6} class, which in
% turn was based upon other sources.  In order to comply with 7th
% Edition criteria, certain changes had to be made to update the
% \textsf{apa6} code.
%
% \subsection{Section Headings}
% Most prominently, the formatting of section headings had to be
% altered.  The 7th Edition specifies a more straightforward series of
% heading levels than previous editions did.  Briefly, the top-level
% heading is now boldfaced and centered, upper- and lower-case, no
% matter now many levels of heading are in the document; other heading
% levels have similar specifications.  The \textsf{apa7} class
% utilizes code (with permission) from the \textsf{apa6} class to
% comply with all of these specifications.
%
% Because the 7th Edition requires figure captions to be printed on
% the same page as their respective figures, there are no more Figure
% Captions pages.
%
% \subsection{Author Note}
% According to 7th Edition guidelines, Author Notes are now placed "in the bottom half of the title page" on the
% title page of manuscripts per the style guide.
%
% \section{Features}
% The featured implemented in the |apa6| have been included in the |apa7| class for the 7th Edition of the
% \textit{Manual}.
%
% \subsection{Masked References}
% \label{sec:masking}
% When manuscripts are sent out for review, they customarily must have
% all identifying information stripped so that reviewers do not know
% who the author of the manuscript is.  The new
% \DescribeMacro{mask}|mask| option suppresses the output of the
% author's name and affiliation, the author note, and any references
% that are marked as being the author's own.
%
% \begin{table}[htbp]
%   \caption{Supported masking commands}
%   \label{tab:CitationCommandsTable}
%   \begin{tabular}{@{}llll@{}}                                           \toprule
%     Unmasked Result    &  \multicolumn{3}{c}{Masking Commands}      \\ \cmidrule(r){2-4}
%                        &  \textsf{apacite}   &  \textsf{natbib}     &  \textsf{biblatex}  \\ \midrule
%     (van Dijk, 2001)  &  |\maskcite|        &  |\maskcitep|        &  |\maskparencite|   \\
%     (Van Dijk, 2001)  &                     &  |\maskCitep|        &  |\maskParencite|   \\
%     van Dijk, 2001    &  |\maskciteNP|      &  |\maskcitealp|      &  |\maskcite|        \\
%     Van Dijk, 2001    &                     &  |\maskCitealp|      &  |\maskCite|        \\
%     van Dijk (2001)   &  |\maskciteA|       &  |\maskcitet|        &  |\masktextcite|    \\
%     Van Dijk (2001)   &                     &  |\maskCitet|        &  |\maskTextcite|    \\
%     van Dijk          &  |\maskciteauthor|  &  |\maskciteauthor|   &  |\maskciteauthor|  \\
%     Van Dijk          &                     &  |\maskCiteauthor|   &  |\maskCiteauthor|  \\
%     (2001)             &  |\maskciteyear|    &  |\maskciteyearpar|  &                     \\
%     2001               &  |\maskciteyearNP|  &  |\maskciteyear|     &  |\maskciteyear|    \\ \bottomrule
%   \end{tabular}
% \end{table}
%
% To specify a reference that is to be masked with the |mask| option,
% simply prepend |mask| to the desired citation command (e.g.,
% |\maskcite| instead of |\cite|).  Supported bibliography packages
% are \textsf{biblatex}, \textsf{apacite}, and \textsf{natbib}.  The
% specific masking commands are shown in Table
% \ref{tab:CitationCommandsTable}.  Prenotes and postnotes for
% citations are supported with both \textsf{biblatex} and
% \textsf{natbib}; however, they are not supported with the
% \textsf{apacite} package.
%
% To mask citations that have been prepended with |mask| (and the
% corresponding bibliography entries), the name of the bibliography
% package must also be passed in as an option on the |\documentclass|
% line, in addition to the |mask| option.  For example,
% |\documentclass[jou,biblatex,mask]{apa7}|. (But as explained later,
% the |biblatex| option is the default and is therefore not required,
% even with the |mask| option.)
%
% \textbf{Be warned}, however, that if you have previously generated
% .bbl, etc., files without the |mask| option applied, those files
% must be deleted or re-written before \LaTeX-ing with the |mask|
% option---otherwise, you will see the to-be-masked entries showing up
% in the bibliography.
%
% Masked citations are replaced with the text, \textit{(2 citations
% removed for masked review)} (in the case of two masked citations).
% The corresponding entries in the References section are also
% suppressed.
%
% There is no need to revise the masked citations when removing the
% |mask| option for final production.  The citations that were
% previously masked will not be masked in the absence of the |mask|
% option.
%
% The |mask| command does not apply to |stu| type documents.
%
% \subsection{Repositioned Floats}
% When revising and proofreading a manuscript, it is most helpful to
% have the tables and figures readily available (rather than turning
% most of the way to the end of the manuscript to access them).  The
% \DescribeMacro{floatsintext}|floatsintext| option (specified in the
% |\documentclass| line) will integrate tables and figures
% approximately where they are mentioned in the text.  This is
% available only with the |man| and |stu| options, of course, because the other
% formats already have floats integrated with the text.
%
% \subsection{User-defined Font Size}
% \DescribeMacro{10pt} \DescribeMacro{11pt} \DescribeMacro{12pt} Users
% can now select from the font-size options available in standard
% \LaTeX\ (10pt, 11pt, 12pt) by including the appropriate option
% (e.g., |10pt|) in the |\documentclass| line.  The default font size
% is 10pt for |jou| mode, 11pt for |doc| mode, and 12pt for |man|
% mode.
%
% \subsection{Watermark}
% If desired, a ``DRAFT'' watermark can be placed on either the first
% pages or all pages of the document with the
% \DescribeMacro{draftfirst}|draftfirst| and
% \DescribeMacro{draftall}|draftall| options, respectively.  The text,
% font size, angle, and lightness of the text can all be modified
% using commands explained in the \textsf{draftwatermark}
% documentation.  Also loads the \textsf{everypage} package.
%
% \subsection{Flexible Bibliographies}
% The \textsf{apa7} class supports three bibliography packages:
% \textsf{biblatex}, \textsf{apacite}, and \textsf{natbib}. \\ \\
% At this time, only \textsf{biblatex} is updated to support the 7th edition of APA style for citations and references.
%
% \subsubsection{\textsf{biblatex}}
% \label{sec:biblatexinfo}
% \textit{\textbf{Note:} This is the recommended method for 7th edition citations and references.} \\
% Biblatex is the most APA-compliant bibiliography package.  The
% \DescribeMacro{biblatex}|biblatex| option directs \textsf{apa7} to
% load the \textsf{biblatex} package with the following option:
% style=apa; however, the |\addbibresource| command will need
% to be specified by the user within the document preamble.  In the
% absence of bibliographic options, \textsf{apa7} will mask references
% using \textsf{biblatex} commands if the |mask| option is specified.
% In other words, the only reason to use the |biblatex| option is to
% have \textsf{apa7} automatically load the \textsf{biblatex} package.
%
% \subsubsection{\textsf{apacite}}
% \label{sec:apaciteinfo}
% \textit{\textbf{Note:} at this time, this will produce citations in 6th edition format.} \\
% The \textsf{apacite} package is loaded if the
% \DescribeMacro{apacite}|apacite| option is specified.  This option
% also informs \textsf{apa7} to mask references using \textsf{apacite}
% commands if the |mask| option is specified.  \textbf{The
% \textsf{apacite} package is \textit{not} loaded by default.}
%
% \subsubsection{\textsf{natbib}}
% \label{sec:natbibinfo}
% \textit{\textbf{Note:} at this time, this will produce citations in 6th edition format.} \\
% Specifying the \DescribeMacro{natbib}|natbib| option implicitly
% loads \textsf{apacite} and \textsf{natbib} and directs
% \textsf{apa7} to mask references using \textsf{natbib} commands if
% the |mask| option is specified.
%
% \subsection{Keywords}
% \label{sec:keywords}
% Many journals (including APA journals) request authors to provide
% keywords for their manuscripts to facilitate electronic indexing.
% \textsf{apa6} introduced the \DescribeMacro{\keywords}|\keywords|
% command.  If provided, keywords will be displayed on a line beneath
% the abstract.  For languages other than English, the `Keywords'
% label can be localized by modifying the appropriate configuration
% file (in the `config' folder of the \textsf{apa7} installation).
%
% \subsection{Float Placement}
% The 7th Edition (as did the 6th Edition) requires
% that tables and figures (in that order) be placed after the
% references but before the appendices.  This creates a bit of a
% conundrum as to what should happen with tables or figures that are
% ultimately typeset within an appendix.  The choices we are left with
% are to place appendix floats (a) along with the floats from the main
% part of the manuscript, which would mean that appendix floats appear
% prior to the point at which they are mentioned; (b) within the
% appendices themselves, which is not consistent with how floats in
% the main part of the manuscript are handled; or (c) in a separate
% float section that follows the appendices, which results in two
% float sections.  Obviously none of these choices is satisfactory, so
% Brian Beitzel (author of |apa6| package) posed the question to APA's Style Expert.  The style expert responded that at
% least for APA's journals ``it doesn't matter whether appendix tables
% are submitted with text tables or separately, as long as they are
% numbered correctly (e.g., Table A1, Table B1, etc.)''  (J.\
% Hume-Pratuch, personal communication, June 15, 2011).  Therefore,
% \textsf{apa7} takes the same approach as \textsf{apa6}: the most straightforward approach that includes
% all appendix floats within the body of the relevant appendix.  This
% also has the advantage of making appendices more readable.
%
% \section{Dependencies}
% \textsf{apa7} automatically loads the following packages.  If these
% packages are not already installed, producing the first
% \textsf{apa7} document could take a few minutes while these packages
% are downloaded and installed.
%
% \begin{itemize}
% \item \textsf{apacite}: bibliography package; used only if the
%   |apacite| or |natbib| option has been specified
% \item \textsf{biblatex}: bibliographic package; used only if the
%   |biblatex| option has been specified
% \item \textsf{booktabs}:  formats tables that are much more
%   attractive than the standard \LaTeX\ tables.
% \item \textsf{caption}: formats table and figure captions
% \item \textsf{draftwatermark}: includes a ``DRAFT'' watermark; used
%   only if the |draftfirst| or |draftall| options are specified
%   (\textsf{draftwatermark} automatically loads \textsf{everypage}
% \item \textsf{endfloat}: handles placing tables and figures at the
%   end of a manuscript; used only with the |man| and |stu| options
% \item \textsf{etoolbox}: provides low-level hooks needed to detect
%   user-loaded packages
% \item \textsf{fancyhdr}: formats page headers
% \item \textsf{float}: handles floats placed within text; used only
%   when the |man| and |floatsintext| options are both specified
%   \textsf{biblatex} package (when loaded by the user in the document
%   preamble) in time to set the ``References'' heading to
%   non-boldface; used only if no bibliographic options have been
%   specified
% \item \textsf{geometry}: formats margins
% \item \textsf{graphicx}: allows inclusion of figures
% \item \textsf{hyperref}: needed for ORCID link
% \item \textsf{longtable}: formats tables that exceed one page in
%   length; loaded only if the |longtable| option is specified.  Also
%   loads the \textsf{array} package.
% \item \textsf{lmodern}: needed for proper default text size for
%   draft watermark, so used on all documents; may be suppressed from
%   loading with option |nolmodern|.
% \item \textsf{scalerel}: needed for ORCID logo sizing
% \item \textsf{substr}: counts masked references; used only if the
%   |mask| option is specified
% \item \textsf{threeparttable}: produces nicely formatted table notes
%   that comply with APA style.  See |longsample.tex| (in the
%   ``samples'' subfolder of the \textsf{apa7} installation;
%   specifically, Appendix B in that document) for how to set up a
%   table with notes.
% \item \textsf{tikz}: needed for ORCID logo
% \item \textsf{times}: for math definitions
% \item \textsf{tracklang}: needed to detect requested languages
% \item \textsf{translations}: for translating strings
% \end{itemize}
%
%
% \section{Examples}
%
% Sample documents are included with this class; look in the
% ``samples'' subfolder of your installation. The source and output for
% |shortsample.tex| are reproduced on the following pages for quick
% reference.
%
% \newpage
%
% \subsection{shortsample.tex}
% \iffalse
%<*verb>
% \fi
 \begin{verbnobox}[\hspace{-3cm}]
 \documentclass[jou]{apa7}

 \usepackage[hungarian,british]{babel}

 \usepackage{csquotes}
 \usepackage[style=apa,backend=biber]{biblatex}
 \addbibresource{bibliography.bib}

 \title{Sample APA-Style Document Using the \textsf{apa7} Package}

 \authorsnames{Daniel A. Weiss}
 \authorsaffiliations{Department of Psychology, A University Somewhere}

 \leftheader{Weiss}

 \abstract{This demonstration paper uses the \textsf{apa7} \LaTeX\
   class to format the document in compliance with the 7th Edition of
   the American Psychological Assocation's \textit{Publication Manual.}
   The references are managed using \textsf{biblatex}.}

 \keywords{APA style, demonstration}

 \authornote{\addORCIDlink{Daniel A. Weiss}{0000-0000-0000-0000}, Department of Psychology, A University Somewhere.

 Changes of Affiliations and other info here to describe what needs to be said}

 \begin{document}
 \maketitle
 We begin with \textcite{Shotton1989}.  We can also cite this work in

 
 Note the use of five heading levels throughout this demonstration
 Method section.

 \section{Method}
 \subsection{Participants}
 We had a lot of people in this study, show in figure
 \ref{fig:figure1}.

 \begin{figure}[!htp]
  \caption{This is my figure caption.}
  \includegraphics[bb=0in 0in 2.5in 2.5in, height=1.4in, width=1.4in]{Figure1.pdf}
  \foreignlanguage{hungarian}{\figurenote{This is an awesome figure with a Hungarian note label.}}
  \label{fig:Figure1}
 \end{figure}

 \subsection{Materials}
 Several materials were used for this project.  Some of them were
 already created for prior research.

 \subsubsection{Paper-and-Pencil Instrument}
 We used an instrument that we found to be highly successful.

 \paragraph{Reliability}
 The reliability of this instrument is extraordinary.

 \paragraph{Validity}
 We now discuss the validity of our instrument.

 \subparagraph{Face validity} The face validity is exceptionally
 strong.  Everyone should be impressed.

 \subparagraph{Construct validity} Also very strong.

 \subsection{Design}
 This section describes the study's design.

 \subsection{Procedure}
 The procedure was fairly straightforward, yet required
 attention to detail.

 \section{Results}
 Table \ref{tab:ComplexTable} contains some sample data.  Our
 statistical prowess in analyzing these data is unmatched.

 \begin{table}[htbp]
   \vspace*{2em}
   \begin{threeparttable}
     \caption{A Complex Table}
     \label{tab:ComplexTable}
     \begin{tabular}{@{}lrrr@{}}         \toprule
     Distribution type  & \multicolumn{2}{l}{Percentage of} & Total number   \\
                        & \multicolumn{2}{l}{targets with}  & of trials per  \\
                        & \multicolumn{2}{l}{segment in}    & participant    \\ \cmidrule(r){2-3}
                                     &  Onset  &  Coda            &          \\ \midrule
     Categorical -- onset\tabfnm{a}  &    100  &     0            &  196     \\
     Probabilistic                   &     80  &    20\tabfnm{*}  &  200     \\
     Categorical -- coda\tabfnm{b}   &      0  &   100\tabfnm{*}  &  196     \\ \midrule
     \end{tabular}
     \tablenote{All data are approximate.
 
             \tabfnt{a}Categorical may be onset.
             \tabfnt{b}Categorical may also be coda.
 
             \tabfnt{*}\textit{p} < .05.
             \tabfnt{**}\textit{p} < .01.
          }
   \end{threeparttable}
 \end{table}

 \section{Discussion}
 This is a lengthy and erudite discussion.  It demonstrates amazing
 skill in interpreting the results for the masses.

 \printbibliography

 \end{document}
 \end{verbnobox}
% \iffalse
%</verb>
% \fi
%
% \subsection{shortsample.pdf}
% \IfFileExists{samples/shortsample.pdf}%
%   {\includegraphics[scale=0.8, trim = 1in 1.5in 1in 0.8in]{samples/shortsample.pdf}}%
%   {PRINTED DOCUMENT WILL BE DISPLAYED HERE \newpage}
%
%
% \StopEventually{^^A
%   \PrintChanges
%   \PrintIndex
% }
%
%    \begin{macrocode}
%<*class>
%    \end{macrocode}
%
%\begin{macro}{apa7.cls}
%    \begin{macrocode}
% % Below is the license statement from the apa class, upon which this
% % apa7 class was built
% %%%%%%%%%%%%%%%%%%%%%%%%%%%%%%%%%%%%%%%%%%%%%%%%%%%%%%%%%%%%%%%%%%%%%%%%%%%%%%
%
% apa6 version 2.33
%
% apa6 - A LaTeX class for formatting documents in compliance with the
% American Psychological Association's Publication Manual, 6th edition
%
% Copyright (C) 2011-2017 by Brian D. Beitzel <brian at beitzel.com>
%
% This work may be distributed and/or modified under the
% conditions of the LaTeX Project Public License (LPPL), either
% version 1.3c of this license or (at your option) any later
% version.  The latest version of this license is in the file:
%
% http://www.latex-project.org/lppl.txt
%
% Users may freely modify these files without permission, as long as the
% copyright line and this statement are maintained intact.
%
% This work is not endorsed by, affiliated with, or probably even known
% by, the American Psychological Association.
%
% %%%%%%%%%%%%%%%%%%%%%%%%%%%%%%%%%%%%%%%%%%%%%%%%%%%%%%%%%%%%%%%%%%%%%%%%%%%%%%

\DeclareOption{man}{%
  \def\def@man{\@manmode}
}

\DeclareOption{stu}{%
  \def\def@stu{\@stumode}
  \def\def@man{\@manmode}
}

\DeclareOption{jou}{%
  \def\def@jou{\@joumode}
}

\DeclareOption{doc}{%
  \def\def@doc{\@docmode}
}

\DeclareOption{babel}{%
  \def\def@babel{\@babel}
}

\DeclareOption{notimes}{%
  \@ifundefined{def@jou}{}{\def\def@notimes{\@notimes}}
}

\DeclareOption{notxfonts}{% -- thp 2005/07/23
  \@ifundefined{def@jou}{}{\def\def@notxfonts{\@notxfonts}}
}

\DeclareOption{nosf}{%
  \@ifundefined{def@man}{}{\def\def@nosf{\@nosf}}
}

\DeclareOption{fignum}{%
  \@ifundefined{def@man}{}{\def\fig@num{\relax}}
}

% \DeclareOption{floatmark}{%
%   \@ifundefined{def@man}{}{\def\float@mark{\relax}}
% }

\DeclareOption{longtable}{%
  \def\long@table{\relax}
}

\DeclareOption{tt}{%
  \@ifundefined{def@man}{}{\def\tt@family{\relax}}
}

\DeclareOption{helv}{%
  \@ifundefined{def@man}{}{\def\helv@family{\relax}}
}

\DeclareOption{notab}{\def\no@tab{\relax}}

\DeclareOption{nobf}{\def\no@bf@title{\relax}}

\DeclareOption{nolmodern}{%
  \def\def@nolmodern{\@nolmodernmode}
}

\DeclareOption{nofontenc}{%
  \def\def@nofontenc{\@nofontencmode}
}

% remove extra space around headings, etc., in man mode
\DeclareOption{noextraspace}{%
  \ClassError{apa7}{Option removed: no longer uses as of version  2.15}{Spacing updated to more closely match APA 7th Edition Manual}
}

% suppress the title at the introduction (in case a different title is desired)
\DeclareOption{donotrepeattitle}{%
  \def\def@donotrepeattitle{\@donotrepeattitlemode}
}

% in order to have tables and figure inside the text, Guillaume Jourjon 19/10/10
\DeclareOption{floatsintext}{%
  \def\def@floatsintext{\@floatsintext}
}

% allow A4 paper size
\DeclareOption{a4paper}{%
  \def\def@aFourPaper{\@aFourPapermode}
}

\DeclareOption{apacite}{% BDB
  \def\def@apacite{\@apacitemode}
}

\DeclareOption{natbib}{% BDB
  \def\def@natbib{\@natbibmode}
}

\DeclareOption{biblatex}{% BDB
  \def\def@biblatex{\@biblatexmode}
}

\DeclareOption{draftfirst}{% BDB
  \def\def@draftfirst{\@draftfirstmode}
}

\DeclareOption{draftall}{% BDB
  \def\def@draftall{\@draftallmode}
}

\DeclareOption{mask}{\def\apaSeven@maskauthoridentity{\relax}}  % BDB

% declare the font-size commands
\newcommand\apaSeven@ptsize{}
\newcommand\apaSeven@noptsize{}
\DeclareOption{10pt}{\renewcommand\apaSeven@ptsize{10pt}}
\DeclareOption{11pt}{\renewcommand\apaSeven@ptsize{11pt}}
\DeclareOption{12pt}{\renewcommand\apaSeven@ptsize{12pt}}


\DeclareOption*{\PassOptionsToClass{\CurrentOption}{article}}

\ProcessOptions\relax

% make jou the default if no other options have been specified
\@ifundefined{def@man}{%
 \@ifundefined{def@doc}{%
  \@ifundefined{def@jou}{%
   \def\def@jou{\@joumode}
   \ClassInfo{apa7}{Using default mode (jou)}
   %\def\def@man{\@manmode}
   %\def\def@doc{\@docmode}
  }{}
 }{}
}{}

% pass the default font-size to the article class
\@ifundefined{def@man}{%
    \@ifundefined{def@jou}{%
        \@ifundefined{def@doc}{%
        }{% doc
          \ifx\apaSeven@ptsize\apaSeven@noptsize
            \LoadClass[11pt]{article} % default for doc is 11pt
          \else
            \LoadClass[\apaSeven@ptsize]{article}
          \fi
        }
    }{% jou
      \ifx\apaSeven@ptsize\apaSeven@noptsize
        \LoadClass[10pt,twoside]{article} % default for jou is 10pt
      \else
        \LoadClass[\apaSeven@ptsize,twoside]{article}
      \fi
    }
}{% man
  \ifx\apaSeven@ptsize\apaSeven@noptsize
    \LoadClass[12pt,twoside]{article} % default for man is 12pt
  \else
    \LoadClass[\apaSeven@ptsize]{article}
  \fi
}

\usepackage{etoolbox}

% determine the bibliography package, and load it if selected through the appropriate \documentclass option
\@ifundefined{def@apacite}{%
  \@ifundefined{def@natbib}{%
    \@ifundefined{def@biblatex}{%
      \def\def@biblatex{\@biblatexmode}%  the default bibliography package is biblatex
      \AtEndPreamble{%
        \@ifpackageloaded{biblatex}{%  the user has loaded biblatex
          \@ifundefined{def@man}{%
            \defbibheading{bibliography}[\refname]{\section*{\normalfont\textbf{#1}}}%
          }{%
            \defbibheading{bibliography}[\refname]{\clearpage\section*{\normalfont\textbf{#1}}}%
          }
        }{}
      }
      \ClassInfo{apa7}{No bibliography package was specified; defaulting to (but not loading) Biblatex}
    }{%
      \def\def@biblatex{\@biblatexmode}%  the selected bibliography package is Biblatex
      \RequirePackage[style=apa,backend=biber]{biblatex}
      \@ifundefined{def@man}{%
        \defbibheading{bibliography}[\refname]{\section*{\normalfont\textbf{#1}}}%
      }{%
        \defbibheading{bibliography}[\refname]{\clearpage\section*{\normalfont\textbf{#1}}}%
      }
      \ClassInfo{apa7}{The selected bibliography package, biblatex, has been loaded}
    }
  }{%
    \def\def@natbib{\@natbibmode}%  the selected bibliography package is natbib (with apacite)
    \@ifundefined{def@man}{%         -- thp 2005/07/23
      \RequirePackage[natbibapa]{apacite}[2012/02/25]}
      {\RequirePackage[natbibapa,bibnewpage]{apacite}[2012/02/25]}
    \ClassInfo{apa7}{The selected bibliography package, apacite and
      natbib, have been loaded}
  }
}{%
  \def\def@apacite{\@apacitemode}%  the selected bibliography package is apacite
  \@ifundefined{def@man}{%         -- thp 2005/07/23
    \RequirePackage{apacite}[2005/06/08]}
    {\RequirePackage[bibnewpage]{apacite}[2005/06/08]}
  \ClassInfo{apa7}{The selected bibliography package, apacite, has been loaded}
}

% load the lmodern package (needed for proper watermark font size)
\@ifundefined{def@nolmodern}{%
  \RequirePackage{lmodern}
}{}

% load the fontenc package (needed for proper hyphenation of accented characters)
\@ifundefined{def@nofontenc}{%
  \RequirePackage[T1]{fontenc}
}{}

% place a watermark on any pages indicated
\@ifundefined{def@draftall}{%
  \@ifundefined{def@draftfirst}{}{%
    \RequirePackage[firstpage]{draftwatermark}
    \SetWatermarkText{DRAFT}
  }
}{%
% a draft watermark has been requested for all pages
  \RequirePackage{draftwatermark}
  \SetWatermarkText{DRAFT}
}

%\@ifundefined{def@noapacite}%       -- Philip Kime 2008/12/03
%  {\ClassInfo{apa7}{Using BibTeX with apacite for citations and references}}%
%  {\ClassInfo{apa7}{apacite not loaded -- User may load BibLaTeX}}


%==== Conditional compilation macros depending on mode ====

\long\def\ifapamodeman#1#2{\@ifundefined{def@man}{#2}{#1}}
\long\def\ifapamodejou#1#2{\@ifundefined{def@jou}{#2}{#1}}
\long\def\ifapamodedoc#1#2{\@ifundefined{def@doc}{#2}{#1}}
\long\def\ifapamode#1#2#3{%
 \@ifundefined{def@man}{%
   \@ifundefined{def@jou}{%
    \@ifundefined{def@doc}{\ClassError{apa7}{Undefined mode state!}}{#3}%
   }{#2}%
  }{#1}%
}

%=================================================================
\@ifundefined{def@man}{}{%
%
% endnotes and endfloat loaded later
% longtable must be loaded before endfloat
%
\@ifundefined{long@table}{}{%
 \RequirePackage{array}
 \RequirePackage{longtable}
%
% longtable messes with captions too, we need to intervene to match style
% NOTE: the caption package makes the following redundant
%\def\LT@makecaption#1#2#3{%
%  \LT@mcol\LT@cols c{\hbox to\z@{\hss\parbox{\linewidth}{%
%      \vskip 10pt%
%      #1{#2}%
%      \let\BBAB\table@BBAB%   -- thp 2005/07/23
%      \par\emph{#3}%
%      \let\BBAB\normal@BBAB%  -- thp 2005/07/23
%    \endgraf\vskip\baselineskip}%
%  \hss}}}
}% END of loading longtable
%
% Switch font family for manuscript if required
%
\@ifundefined{tt@family}{}{%
% A typewriter font can be used as in the APA manual example. Since tt has
% no bold (or bold-italic), the cmss font metrics are loaded instead
% \DeclareFontShape{OT1}{cmtt}{bx}{n}{ <-> sub * cmss/bx/n }{} % this fails
 \DeclareFontShape{OT1}{cmtt}{bx}{n}{ <-> cmssbx10 }{}  % probably not the
 \DeclareFontShape{OT1}{cmtt}{bx}{it}{ <-> cmssbxo10}{} % right way to do it
 \renewcommand{\familydefault}{cmtt}
 }
\@ifundefined{helv@family}{}{%
 \renewcommand{\familydefault}{phv}}
}


%
% For jou mode, load times (and related math fonts) if available
% First try txfonts because they are supposed to have better math definitions
%
\@ifundefined{def@jou}{}{%
 \@ifundefined{def@notimes}{%
  \newif\iftxfonts          % -- thp 2005/07/23
  \txfontsfalse             % added checks for txfonts because they may be undesirable
                            % for example, there are no Greek txfonts but there are times
  \IfFileExists{txfonts.sty}{\@ifundefined{def@notxfonts}{\txfontstrue}{}}{}
   \iftxfonts%
    \RequirePackage{txfonts}%
    \typeout{apa7.cls: Using txfonts}% Changed from Warning -- thp 2005/12/28
    %%%
    % According to Erik Meijer, txfonts causes problems if amsmath is loaded later
    % (i.e., via \usepackage by the user); instead of providing yet another option
    % to load amsmath by apa.cls, we adopt Erik's suggestion to undefine temporarily
    % the offending macros -- thp 2005/12/28
    \let\tempiint\iint\let\iint\undefined
    \let\tempiiint\iiint\let\iiint\undefined
    \let\tempiiiint\iiiint\let\iiiint\undefined
    \let\tempidotsint\idotsint\let\idotsint\undefined
    \let\tempopenbox\openbox\let\openbox\undefined
    \AtBeginDocument{%
     \let\iint\tempiint\let\tempiint\undefined
     \let\iiint\tempiiint\let\tempiiint\undefined
     \let\iiiint\tempiiiint\let\tempiiiint\undefined
     \let\idotsint\tempidotsint\let\tempidotsint\undefined
     \let\openbox\tempopenbox\let\tempopenbox\undefined
    }
    %%% end of taking care of txfonts problems
   \else%
    % if txfonts are not available/desirable, try pslatextimes/mathptm
    \IfFileExists{pslatex.sty}
     {\RequirePackage{pslatex}}
     % if pslatex is not available, try times/mathptm
     {\RequirePackage{times}
      \IfFileExists{mathptm.sty}{\RequirePackage{mathptm}}{}}%
   \fi% txfonts not available/desirable
 }{}% def@notimes
}% def@jou


% choose a paper size and set the page margins
\@ifundefined{def@aFourPaper}{
  \RequirePackage[top=1in, bottom=1in, left=1in, right=1in]{geometry}  
}{
  \RequirePackage[top=1in, bottom=1in, left=1in, right=1in,a4paper]{geometry}  
}

\RequirePackage{graphicx}  % this is for including graphics

\RequirePackage{scalerel,tikz,hyperref} % this is included for ORCID icon
\usetikzlibrary{svg.path}
\definecolor{orcidlogocol}{HTML}{A6CE39}
\tikzset{
  orcidlogo/.pic={
    \fill[orcidlogocol] svg{M256,128c0,70.7-57.3,128-128,128C57.3,256,0,198.7,0,128C0,57.3,57.3,0,128,0C198.7,0,256,57.3,256,128z};
    \fill[white] svg{M86.3,186.2H70.9V79.1h15.4v48.4V186.2z}
                 svg{M108.9,79.1h41.6c39.6,0,57,28.3,57,53.6c0,27.5-21.5,53.6-56.8,53.6h-41.8V79.1z M124.3,172.4h24.5c34.9,0,42.9-26.5,42.9-39.7c0-21.5-13.7-39.7-43.7-39.7h-23.7V172.4z}
                 svg{M88.7,56.8c0,5.5-4.5,10.1-10.1,10.1c-5.6,0-10.1-4.6-10.1-10.1c0-5.6,4.5-10.1,10.1-10.1C84.2,46.7,88.7,51.3,88.7,56.8z};
  }
}
\newcommand{\addORCIDlink}[2]{#1 \href{https://orcid.org/#2}{{\mbox{\scalerel*{
\begin{tikzpicture}[yscale=-1,transform shape]
\pic{orcidlogo};
\end{tikzpicture}
}{|}}} https://orcid.org/#2}}

\RequirePackage{booktabs}  % this is for nice-looking tables
\setlength{\abovetopsep}{0pt}  % set the distance between the table title and the table toprule
\setlength{\belowbottomsep}{0pt}  % set the distance between the table bottomrule and any notes

\RequirePackage[para,flushleft]{threeparttable}  % this is for nice-looking table footnotes, etc.
% need to redefine the tablenotes so they hug the bottom line of the table a little closer
% also, the manuscript format needs to have \raggedright
\@ifundefined{def@man}{% BDB
  \def\TPT@doparanotes{\par\vspace{-.5\baselineskip}% BDB
     \prevdepth\z@ \TPT@hsize
     \TPTnoteSettings
     \parindent\z@ \pretolerance 8
     \linepenalty 200
     \renewcommand\item[1][]{\relax\ifhmode \begingroup
         \unskip
         \advance\hsize 10em % \hsize is scratch register, based on real hsize
         \penalty -45 \hskip\z@\@plus\hsize \penalty-19
         \hskip .15\hsize \penalty 9999 \hskip-.15\hsize
         \hskip .01\hsize\@plus-\hsize\@minus.01\hsize
         \hskip 1em\@plus .3em
        \endgroup\fi
        \tnote{##1}\,\ignorespaces}%
     \let\TPToverlap\relax
     \def\endtablenotes{\par}%
  }
}{%
  \def\TPT@doparanotes{\par\vspace{-.4\baselineskip}% BDB
     \prevdepth\z@ \TPT@hsize
     \TPTnoteSettings
     \raggedright
     \parindent\z@ \pretolerance 8
     \linepenalty 200
     \renewcommand\item[1][]{\relax\ifhmode \begingroup
         \unskip
         \advance\hsize 10em % \hsize is scratch register, based on real hsize
         \penalty -45 \hskip\z@\@plus\hsize \penalty-19
         \hskip .15\hsize \penalty 9999 \hskip-.15\hsize
         \hskip .01\hsize\@plus-\hsize\@minus.01\hsize
         \hskip 1em\@plus .3em
        \endgroup\fi
        \tnote{##1}\,\ignorespaces}%
     \let\TPToverlap\relax
     \def\endtablenotes{\par}%
  }
}

%========= Language portability parameters ===========

% Language-specific definition of strings is handled externally -- thp 2005/12/30
% Based on handling of the same issue in apacite.sty

\RequirePackage{translations}

% First define the default American English macros in case no external file is present or needed
\DeclareTranslationFallback{prelastauthorsep}{,}
\DeclareTranslationFallback{lastauthorseparator}{and}
\DeclareTranslationFallback{acksname}{Author Note}
\DeclareTranslationFallback{keywordname}{Keywords}
\DeclareTranslationFallback{notesname}{Footnotes}
\DeclareTranslationFallback{notelabel}{Note}

\AtEndPreamble{%
  \RequirePackage{tracklang}%
  \TrackLangShowVerbosetrue

  \ForEachTrackedDialect{\thisdialect}{%
    \LoadDictionaryFor{\thisdialect}{apa7}%
  }
}

%========== Load babel if required ===================

\@ifundefined{def@babel}{}{% -- thp 2005/07/23
 \RequirePackage{babel}    % -- thp 2005/07/23, removed options 2005/12/28
}

\@ifundefined{def@biblatex}{}{% BDB

  % we are using biblatex

    %%%%%%%%%%%% biblatex commands %%%%%%%%%%%%%%%%
    %%
    %%  \cite[e.g.,][p.~11]{vanDijk2001,Ross1987}          =>  e.g., van Dijk, 2001; Ross, 1987, p. 11
    %%  \Cite[e.g.,][p.~11]{vanDijk2001,Ross1987}          =>  e.g., Van Dijk, 2001; Ross, 1987, p. 11
    %%  \parencite[e.g.,][p.~11]{vanDijk2001,Ross1987}     =>  (e.g., van Dijk, 2001; Ross, 1987, p. 11)
    %%  \Parencite[e.g.,][p.~11]{vanDijk2001,Ross1987}     =>  (e.g., Van Dijk, 2001; Ross, 1987, p. 11)
    %%  \textcite[e.g.,][p.~11]{vanDijk2001,Ross1987}      =>  e.g., van Dijk (2001); Ross (1987, p. 11)
    %%  \Textcite[e.g.,][p.~11]{vanDijk2001,Ross1987}      =>  e.g., Van Dijk (2001); Ross (1987, p. 11)
    %%  \citeauthor[e.g.,][p.~11]{vanDijk2001,Ross1987}    =>  e.g., van Dijk (2001); Ross (1987, p. 11)
    %%  \Citeauthor[e.g.,][p.~11]{vanDijk2001,Ross1987}    =>  e.g., Van Dijk (2001); Ross (1987, p. 11)
    %%  \citeyear[e.g.,][p.~11]{vanDijk2001}             =>  e.g., 2001, p. 11)
    %%  \footcite[e.g.,][p.~11]{vanDijk2001,Ross1987}      =>  e.g., van Dijk, 2001; Ross, 1987, p. 11. [as footnote]
    %%  \footcitetext[e.g.,][p.~11]{vanDijk2001,Ross1987}  =>  e.g., van Dijk, 2001; Ross, 1987, p. 11. [as footnotetext]
    %%
    %%%%%%%%%%%%%%%%%%%%%%%%%%%%%%%%%%%%%%%%%%%%%%%

    \@ifundefined{apaSeven@maskauthoridentity}%  BDB
      {%  change masked references to unmasked
        \providecommand\maskcite\cite
        \providecommand\maskCite\Cite
        \providecommand\maskparencite\parencite
        \providecommand\maskParencite\Parencite
        \providecommand\masktextcite\textcite
        \providecommand\maskTextcite\Textcite
        \providecommand\maskciteauthor\citeauthor
        \providecommand\maskCiteauthor\Citeauthor
        \providecommand\maskciteyear\citeyear
        \providecommand\maskfootcite\footcite
        \providecommand\maskfootcitetext\footcitetext
      }{%  mask references to author

        \RequirePackage{substr}  % to allow counting of masked references
        \newcounter{maskedRefs}

        % \maskcite
        \newcommand\maskcite{\@ifnextchar[{\maskcite@@also}{\maskcite@@also[]}}
        \newcommand\maskcite@@also{}
        \def\maskcite@@also[#1]{\@ifnextchar[{\maskcite@@@also[#1]}{\maskcite@@@also[][#1]}}

        \def\maskcite@@@also%
            [#1][#2]#3{%
                \setcounter{maskedRefs}{0}%
                \SubStringsToCounter{maskedRefs}{,}{#3}%
                \addtocounter{maskedRefs}{1}%
                \ifnum\value{maskedRefs} = 1%
                \def\apaSeven@masked@refs{(\textit{\themaskedRefs\ citation removed for masked review})}%
                \else%
                \def\apaSeven@masked@refs{(\textit{\themaskedRefs\ citations removed for masked review})}%
                \fi%
                \apaSeven@masked@refs%
        }

        % \maskCite
        \newcommand\maskCite{\@ifnextchar[{\maskCite@@also}{\maskCite@@also[]}}
        \newcommand\maskCite@@also{}
        \def\maskCite@@also[#1]{\@ifnextchar[{\maskCite@@@also[#1]}{\maskCite@@@also[][#1]}}

        \def\maskCite@@@also%
            [#1][#2]#3{%
                \setcounter{maskedRefs}{0}%
                \SubStringsToCounter{maskedRefs}{,}{#3}%
                \addtocounter{maskedRefs}{1}%
                \ifnum\value{maskedRefs} = 1%
                \def\apaSeven@masked@refs{(\textit{\themaskedRefs\ citation removed for masked review})}%
                \else%
                \def\apaSeven@masked@refs{(\textit{\themaskedRefs\ citations removed for masked review})}%
                \fi%
                \apaSeven@masked@refs%
        }

        % \maskparencite
        \newcommand\maskparencite{\@ifnextchar[{\maskparencite@@also}{\maskparencite@@also[]}}
        \newcommand\maskparencite@@also{}
        \def\maskparencite@@also[#1]{\@ifnextchar[{\maskparencite@@@also[#1]}{\maskparencite@@@also[][#1]}}

        \def\maskparencite@@@also%
            [#1][#2]#3{%
                \setcounter{maskedRefs}{0}%
                \SubStringsToCounter{maskedRefs}{,}{#3}%
                \addtocounter{maskedRefs}{1}%
                \ifnum\value{maskedRefs} = 1%
                \def\apaSeven@masked@refs{(\textit{\themaskedRefs\ citation removed for masked review})}%
                \else%
                \def\apaSeven@masked@refs{(\textit{\themaskedRefs\ citations removed for masked review})}%
                \fi%
                \apaSeven@masked@refs%
        }

        % \maskParencite
        \newcommand\maskParencite{\@ifnextchar[{\maskParencite@@also}{\maskParencite@@also[]}}
        \newcommand\maskParencite@@also{}
        \def\maskParencite@@also[#1]{\@ifnextchar[{\maskParencite@@@also[#1]}{\maskParencite@@@also[][#1]}}

        \def\maskParencite@@@also%
            [#1][#2]#3{%
                \setcounter{maskedRefs}{0}%
                \SubStringsToCounter{maskedRefs}{,}{#3}%
                \addtocounter{maskedRefs}{1}%
                \ifnum\value{maskedRefs} = 1%
                \def\apaSeven@masked@refs{(\textit{\themaskedRefs\ citation removed for masked review})}%
                \else%
                \def\apaSeven@masked@refs{(\textit{\themaskedRefs\ citations removed for masked review})}%
                \fi%
                \apaSeven@masked@refs%
        }

        % \maskciteauthor
        \newcommand\maskciteauthor{\@ifnextchar[{\maskciteauthor@@also}{\maskciteauthor@@also[]}}
        \newcommand\maskciteauthor@@also{}
        \def\maskciteauthor@@also[#1]{\@ifnextchar[{\maskciteauthor@@@also[#1]}{\maskciteauthor@@@also[][#1]}}

        \def\maskciteauthor@@@also%
            [#1][#2]#3{%
                \setcounter{maskedRefs}{0}%
                \SubStringsToCounter{maskedRefs}{,}{#3}%
                \addtocounter{maskedRefs}{1}%
                \ifnum\value{maskedRefs} = 1%
                \def\apaSeven@masked@refs{(\textit{\themaskedRefs\ citation removed for masked review})}%
                \else%
                \def\apaSeven@masked@refs{(\textit{\themaskedRefs\ citations removed for masked review})}%
                \fi%
                \apaSeven@masked@refs%
        }

        % \maskCiteauthor
        \newcommand\maskCiteauthor{\@ifnextchar[{\maskCiteauthor@@also}{\maskCiteauthor@@also[]}}
        \newcommand\maskCiteauthor@@also{}
        \def\maskCiteauthor@@also[#1]{\@ifnextchar[{\maskCiteauthor@@@also[#1]}{\maskCiteauthor@@@also[][#1]}}

        \def\maskCiteauthor@@@also%
            [#1][#2]#3{%
                \setcounter{maskedRefs}{0}%
                \SubStringsToCounter{maskedRefs}{,}{#3}%
                \addtocounter{maskedRefs}{1}%
                \ifnum\value{maskedRefs} = 1%
                \def\apaSeven@masked@refs{(\textit{\themaskedRefs\ citation removed for masked review})}%
                \else%
                \def\apaSeven@masked@refs{(\textit{\themaskedRefs\ citations removed for masked review})}%
                \fi%
                \apaSeven@masked@refs%
        }

        % \maskciteyear
        \newcommand\maskciteyear{\@ifnextchar[{\maskciteyear@@also}{\maskciteyear@@also[]}}
        \newcommand\maskciteyear@@also{}
        \def\maskciteyear@@also[#1]{\@ifnextchar[{\maskciteyear@@@also[#1]}{\maskciteyear@@@also[][#1]}}

        \def\maskciteyear@@@also%
            [#1][#2]#3{%
                \setcounter{maskedRefs}{0}%
                \SubStringsToCounter{maskedRefs}{,}{#3}%
                \addtocounter{maskedRefs}{1}%
                \ifnum\value{maskedRefs} = 1%
                \def\apaSeven@masked@refs{(\textit{\themaskedRefs\ citation removed for masked review})}%
                \else%
                \def\apaSeven@masked@refs{(\textit{\themaskedRefs\ citations removed for masked review})}%
                \fi%
                \apaSeven@masked@refs%
        }

        % \maskfootcite
        \newcommand\maskfootcite{\@ifnextchar[{\maskfootcite@@also}{\maskfootcite@@also[]}}
        \newcommand\maskfootcite@@also{}
        \def\maskfootcite@@also[#1]{\@ifnextchar[{\maskfootcite@@@also[#1]}{\maskfootcite@@@also[][#1]}}

        \def\maskfootcite@@@also%
            [#1][#2]#3{%
                \setcounter{maskedRefs}{0}%
                \SubStringsToCounter{maskedRefs}{,}{#3}%
                \addtocounter{maskedRefs}{1}%
                \ifnum\value{maskedRefs} = 1%
                \def\apaSeven@masked@refs{(\textit{\themaskedRefs\ citation removed for masked review})}%
                \else%
                \def\apaSeven@masked@refs{(\textit{\themaskedRefs\ citations removed for masked review})}%
                \fi%
                \apaSeven@masked@refs%
        }

        % \maskfootcitetext
        \newcommand\maskfootcitetext{\@ifnextchar[{\maskfootcitetext@@also}{\maskfootcitetext@@also[]}}
        \newcommand\maskfootcitetext@@also{}
        \def\maskfootcitetext@@also[#1]{\@ifnextchar[{\maskfootcitetext@@@also[#1]}{\maskfootcitetext@@@also[][#1]}}

        \def\maskfootcitetext@@@also%
            [#1][#2]#3{%
                \setcounter{maskedRefs}{0}%
                \SubStringsToCounter{maskedRefs}{,}{#3}%
                \addtocounter{maskedRefs}{1}%
                \ifnum\value{maskedRefs} = 1%
                \def\apaSeven@masked@refs{(\textit{\themaskedRefs\ citation removed for masked review})}%
                \else%
                \def\apaSeven@masked@refs{(\textit{\themaskedRefs\ citations removed for masked review})}%
                \fi%
                \apaSeven@masked@refs%
        }

        % \masktextcite
        \newcommand\masktextcite{\@ifnextchar[{\masktextcite@@also}{\masktextcite@@also[]}}
        \newcommand\masktextcite@@also{}
        \def\masktextcite@@also[#1]{\@ifnextchar[{\masktextcite@@@also[#1]}{\masktextcite@@@also[][#1]}}

        \def\masktextcite@@@also%
            [#1][#2]#3{%
                \setcounter{maskedRefs}{0}%
                \SubStringsToCounter{maskedRefs}{,}{#3}%
                \addtocounter{maskedRefs}{1}%
                \ifnum\value{maskedRefs} = 1%
                \def\apaSeven@masked@refs{(\textit{\themaskedRefs\ citation removed for masked review})}%
                \else%
                \def\apaSeven@masked@refs{(\textit{\themaskedRefs\ citations removed for masked review})}%
                \fi%
                \apaSeven@masked@refs%
        }

        % \maskTextcite
        \newcommand\maskTextcite{\@ifnextchar[{\maskTextcite@@also}{\maskTextcite@@also[]}}
        \newcommand\maskTextcite@@also{}
        \def\maskTextcite@@also[#1]{\@ifnextchar[{\maskTextcite@@@also[#1]}{\maskTextcite@@@also[][#1]}}

        \def\maskTextcite@@@also%
            [#1][#2]#3{%
                \setcounter{maskedRefs}{0}%
                \SubStringsToCounter{maskedRefs}{,}{#3}%
                \addtocounter{maskedRefs}{1}%
                \ifnum\value{maskedRefs} = 1%
                \def\apaSeven@masked@refs{(\textit{\themaskedRefs\ citation removed for masked review})}%
                \else%
                \def\apaSeven@masked@refs{(\textit{\themaskedRefs\ citations removed for masked review})}%
                \fi%
                \apaSeven@masked@refs%
        }

      }

}



%========== APA Section spacing

\newskip\b@level@one@skip   \b@level@one@skip=2.5ex plus 1ex minus .2ex
\newskip\e@level@one@skip   \e@level@one@skip=1.5ex plus .6ex minus .1ex
\newskip\b@level@two@skip   \b@level@two@skip=2.5ex plus 1ex minus .2ex
\newskip\e@level@two@skip   \e@level@two@skip=1.5ex plus .6ex minus .1ex
\newskip\b@level@three@skip \b@level@three@skip=2.0ex plus .8ex minus .2ex
\newskip\e@level@three@skip \e@level@three@skip=1.5ex plus .6ex minus .1ex
\newskip\b@level@four@skip  \b@level@four@skip=0\baselineskip \@plus 0.2ex \@minus 0.2ex
\newskip\e@level@four@skip  \e@level@four@skip=-0.5em
\newskip\b@level@five@skip  \b@level@five@skip=0\baselineskip \@plus 0.2ex \@minus 0.2ex
\newskip\e@level@five@skip  \e@level@five@skip=-0.5em

\ifapamodeman{%
    \b@level@one@skip=0.2\baselineskip \@plus 0.2ex \@minus 0.2ex
    \e@level@one@skip=0.2\baselineskip \@plus .2ex
    \b@level@two@skip=0.2\baselineskip \@plus 0.2ex \@minus 0.2ex
    \e@level@two@skip=0.2\baselineskip \@plus 0.2ex
    \b@level@three@skip=0.2\baselineskip \@plus 0.2ex \@minus 0.2ex
    \e@level@three@skip=0.2\baselineskip \@plus 0.2ex
    \b@level@four@skip=0\baselineskip \@plus 0.2ex \@minus 0.2ex
    \e@level@four@skip=-0.5em
    \b@level@five@skip=0\baselineskip \@plus 0.2ex \@minus 0.2ex
    \e@level@five@skip=-0.5em
}{}

%========== APA Section heading & seriation, adapted from class apa6e

% % Below is the license statement from the apa6e class, from which a
% % small bit of code was used (with permission from Nathaniel Smith)
% % Copyright (C) 2011 by Nathaniel J. Smith <njs@pobox.com>
% %
% % Redistribution and use in source and binary forms, with or without
% % modification, are permitted provided that the following conditions are
% % met:
% %
% %    1. Redistributions of source code must retain the above copyright
% %       notice, this list of conditions and the following disclaimer.
% %
% %    2. Redistributions in binary form must reproduce the above
% %       copyright notice, this list of conditions and the following
% %       disclaimer in the documentation and/or other materials provided
% %       with the distribution.
% %
% % THIS SOFTWARE IS PROVIDED BY THE COPYRIGHT HOLDERS AND CONTRIBUTORS
% % ``AS IS'' AND ANY EXPRESS OR IMPLIED WARRANTIES, INCLUDING, BUT NOT
% % LIMITED TO, THE IMPLIED WARRANTIES OF MERCHANTABILITY AND FITNESS FOR
% % A PARTICULAR PURPOSE ARE DISCLAIMED. IN NO EVENT SHALL THE COPYRIGHT
% % HOLDER OR CONTRIBUTORS BE LIABLE FOR ANY DIRECT, INDIRECT, INCIDENTAL,
% % SPECIAL, EXEMPLARY, OR CONSEQUENTIAL DAMAGES (INCLUDING, BUT NOT
% % LIMITED TO, PROCUREMENT OF SUBSTITUTE GOODS OR SERVICES; LOSS OF USE,
% % DATA, OR PROFITS; OR BUSINESS INTERRUPTION) HOWEVER CAUSED AND ON ANY
% % THEORY OF LIABILITY, WHETHER IN CONTRACT, STRICT LIABILITY, OR TORT
% % (INCLUDING NEGLIGENCE OR OTHERWISE) ARISING IN ANY WAY OUT OF THE USE
% % OF THIS SOFTWARE, EVEN IF ADVISED OF THE POSSIBILITY OF SUCH DAMAGE.
% %

\setcounter{secnumdepth}{0}

\renewcommand{\section}{\@startsection {section}{1}{\z@}%
    {\b@level@one@skip}{\e@level@one@skip}%
    {\centering\normalfont\normalsize\bfseries}}

\renewcommand{\subsection}{\@startsection{subsection}{2}{\z@}%
    {\b@level@two@skip}{\e@level@two@skip}%
    {\normalfont\normalsize\bfseries}}


\newcommand*{\typesectitle}[1]{#1\addperi}

\newcommand*{\addperi}{%
  \relax\ifhmode%
    \ifnum\spacefactor>\@m \else.\fi%
  \fi%
}

\renewcommand{\subsubsection}{\@startsection{subsubsection}{3}{\z@}%
    {\b@level@three@skip}{\e@level@three@skip}%
    {\normalfont\normalsize\bfseries\itshape}}

\renewcommand{\paragraph}{\@startsection{paragraph}{4}{\parindent}%
    {\b@level@four@skip}%
    {\e@level@four@skip}%
    {\normalfont\normalsize\bfseries\typesectitle}}

\renewcommand{\subparagraph}[1]{\@startsection{subparagraph}{5}{1em}%
    {\b@level@five@skip}%
    {\e@level@five@skip}%
    {\normalfont\normalsize\bfseries\itshape\hspace{\parindent}{#1}\textit{\addperi}}{\relax}}

% make the References section non-boldface
\AtBeginDocument{\def\st@rtbibsection{\mspart{\refname}}}%  BDB -- this is for apacite
\AtBeginDocument{\def\bibsection{\mspart{\refname}}}%  BDB -- this is for apacite + natbib
\newcommand{\mspart}{{\ifapamodeman{\clearpage}{}}\@startsection {section}{1}{\z@}%
    {\b@level@one@skip}{\e@level@one@skip}%
    {\centering\normalfont\normalsize}}

% format the table and figure captions
\RequirePackage[singlelinecheck=off]{caption}

\ifapamode{% man
   \DeclareCaptionLabelFormat{tablelabel}{\hspace{-\parindent}\raggedright\textbf{#1 #2}}
    \DeclareCaptionLabelFormat{figurelabel}{\hspace{-\parindent}\raggedright\textbf{#1 #2}}
    \DeclareCaptionTextFormat{tabletext}{\hspace{-\parindent}\raggedright\textit{#1}}
    \DeclareCaptionLabelSeparator{apalabelsep}{\\}
}{% jou
    \DeclareCaptionLabelFormat{tablelabel}{\hspace{-\parindent}\raggedright\textbf{#1 #2}}
    \DeclareCaptionLabelFormat{figurelabel}{\hspace{-\parindent}\raggedright\textbf{#1 #2}}
    \DeclareCaptionTextFormat{tabletext}{\hspace{-\parindent}\textit{#1}}
    \DeclareCaptionLabelSeparator{apalabelsep}{\\[\baselineskip]} 
}{% doc
    \DeclareCaptionLabelFormat{tablelabel}{\hspace{-\parindent}\raggedright\textbf{#1 #2}}
    \DeclareCaptionLabelFormat{figurelabel}{\hspace{-\parindent}\raggedright\textbf{#1 #2}}
    \DeclareCaptionTextFormat{tabletext}{\hspace{-\parindent}\raggedright\textit{#1}}
    \DeclareCaptionLabelSeparator{apalabelsep}{\\[\baselineskip]} 
}
\captionsetup[table]{position=above,skip=0pt,labelformat=tablelabel,labelsep=apalabelsep,textformat=tabletext}
\captionsetup[figure]{position=above,skip=0pt,labelformat=figurelabel,labelsep=apalabelsep,textfont=it}

% space between caption and surrounding document text
% \setlength{\belowcaptionskip}{0pt}


%=============================================================
%
% End of code from apa6e
%
%=============================================================

\newcommand{\figurenote}[1]{\par \small \textit{\GetTranslation{notelabel}.} {#1}}
\newcommand{\tablenote}[1]{\begin{tablenotes}[para,flushleft]
        {\small
            \textit{\GetTranslation{notelabel}.} {#1}
         }
    \end{tablenotes}}

% Code from theapa.sty

% Refer to LaTeX book to modify, if you want, spaces before and after of
%  \begin{...} ... \end{...} or spaces between \item-s.
\newcounter{APAenum}
\newskip\@text@par@indent
\def\APAenumerate{\@text@par@indent\parindent\setbox0\hbox{1. }%
    \list{\arabic{APAenum}.}{\usecounter{APAenum}
    \labelwidth\z@\labelsep\z@\leftmargin\z@\parsep\z@
    \rightmargin\z@\itemsep\z@\topsep\z@\partopsep\z@
    \itemindent\@text@par@indent\advance\itemindent by\wd0
    \def\makelabel##1{\hss\llap{##1 }}}}
\let\endAPAenumerate=\endlist

\def\seriate{\@bsphack\begingroup%
   \setcounter{APAenum}{0}%
   \def\item{\addtocounter{APAenum}{1}(\alph{APAenum})\space}%
   \ignorespaces}
\def\endseriate{\endgroup\@esphack}

\def\APAitemize{\@text@par@indent\parindent\setbox0\hbox{$\bullet$}%
    \list{$\bullet$}{%
    \labelwidth\z@\labelsep.5em\leftmargin\z@\parsep\z@
    \rightmargin\z@\itemsep\z@\topsep\z@\partopsep\z@
    \itemindent\@text@par@indent
    \advance\itemindent by\wd0\advance\itemindent by.5em
    \def\makelabel##1{\hss\llap{##1}}}}
\let\endAPAitemize=\endlist

% \let\enumerate=\APAenumerate
% \let\endenumerate=\endAPAenumerate
% \let\itemize=\APAitemize
% \let\enditemize=\endAPAitemize

%=============================================================
%
% End of code from theapa.sty
%
%=============================================================

%===== apa.cls main declarations for title page contents =====

\long\def\title#1{\long\def\@title{#1}}
\long\def\author#1{\authorsnames{#1}}
\long\def\affiliation#1{\ClassWarning{apa7}{Deprecated: see documentation for \string\authorsaffiliations}\authorsaffiliations{#1}}
\long\def\twoauthors#1#2{\ClassError{apa7}{Deprecated: see documentation for \string\authorsnames}{For multiple authors use \string\authorsnames}}
\long\def\threeauthors#1#2#3{\ClassError{apa7}{Deprecated: see documentation for \string\authorsnames}{For multiple authors use \string\authorsnames}}
\long\def\onetwoauthors#1#2#3{\ClassError{apa7}{Deprecated: see documentation for \string\authorsnames}{For multiple authors use \string\authorsnames}}
\long\def\twooneauthors#1#2#3{\ClassError{apa7}{Deprecated: see documentation for \string\authorsnames}{For multiple authors use \string\authorsnames}}
\long\def\fourauthors#1#2#3#4{\ClassError{apa7}{Deprecated: see documentation for \string\authorsnames}{For multiple authors use \string\authorsnames}}
\long\def\fiveauthors#1#2#3#4#5{\ClassError{apa7}{Deprecated: see documentation for \string\authorsnames}{For multiple authors use \string\authorsnames}}
\long\def\sixauthors#1#2#3#4#5#6{\ClassError{apa7}{Deprecated: see documentation for \string\authorsnames}{For multiple authors use \string\authorsnames}}
\long\def\twoaffiliations#1#2{\ClassError{apa7}{Deprecated: see documentation for \string\authorsaffiliations}{For multiple affiliations use \string\authorsaffiliations}}
\long\def\threeaffiliations#1#2#3{\ClassError{apa7}{Deprecated: see documentation for \string\authorsaffiliations}{For multiple affiliations use \string\authorsaffiliations}}
\long\def\fouraffiliations#1#2#3#4{\ClassError{apa7}{Deprecated: see documentation for \string\authorsaffiliations}{For multiple affiliations use \string\authorsaffiliations}}
\long\def\fiveaffiliations#1#2#3#4#5{\ClassError{apa7}{Deprecated: see documentation for \string\authorsaffiliations}{For multiple affiliations use \string\authorsaffiliations}}
\long\def\sixaffiliations#1#2#3#4#5#6{\ClassError{apa7}{Deprecated: see documentation for \string\authorsaffiliations}{For multiple affiliations use \string\authorsaffiliations}}
\long\def\course#1{\long\def\@course{#1}}
\long\def\professor#1{\long\def\@professor{#1}}
\long\def\duedate#1{\long\def\@duedate{#1}}
\long\def\shorttitle#1{\long\def\@shorttitle{#1}}
\long\def\note#1{\long\def\@note{#1}}
\long\def\abstract#1{\long\def\@abstract{#1}}
\long\def\keywords#1{\long\def\@keywords{#1}}
%\long\def\acknowledgements#1{\long\def\@acks{#1}}
\long\def\authornote#1{\long\def\@acks{#1}}
\def\journal#1{\def\@journal{#1}}
\def\volume#1{\def\@vvolume{#1}}
\def\ccoppy#1{\def\@ccoppy{#1}}
\def\copnum#1{\def\@copnum{#1}}

% Combined code from https://tex.stackexchange.com/questions/297507/how-to-apply-a-do-function-to-two-lists and 
% https://www.dickimaw-books.com/latex/admin/html/exercises/oxfordcomma.shtml

% Create etoolbox lists for multiple authors and affiliations
\newcommand*\listauthors{}
\newcommand*\listsuperscripts{}
\newcommand*\listaffiliations{}

\newcommand*{\authorsnames}[2][]{
  \def\def@multipleauthors{\@multipleauthorsmode} %
  \renewcommand*\listauthors{}
  \renewcommand*\listsuperscripts{}
  \newcounter{NumberOfAuthors}
  \newcounter{NumberOfSuperscripts}
  \forcsvlist{\stepcounter{NumberOfAuthors}\listadd\listauthors}{#2}
  \forcsvlist{\stepcounter{NumberOfSuperscripts}\listadd\listsuperscripts}{#1}
  }
  
  \newcommand*{\authorsaffiliations}[1]{
  \def\def@multipleaffils{\@multipleaffilsmode} %
  \renewcommand*\listaffiliations{}
  \newcounter{NumberOfAffiliations}
  \forcsvlist{\stepcounter{NumberOfAffiliations}\listadd\listaffiliations}{#1}
  }

\makeatletter
  \catcode`\|=3
  
\def\looptwo#1#2{%
    \edef\tmp{\noexpand\xtwo%
      \unexpanded\expandafter{#1}\relax  % no added delimiter here
      \unexpanded\expandafter{#2}\relax  % no added delimiter here
    } \tmp%
  }%
  
\def\xtwo#1|#2\relax#3|#4\relax{%
    \dotwo{#1}{#3}%
    \def\tmp{#2}%
    \ifx\empty\tmp%
      \expandafter\@gobble%
    \else%
      \expandafter\@firstofone%
    \fi%
    {\xtwo#2\relax#4\relax}%
  }%

  \catcode`\|=12
\makeatother

\newcommand*{\dotwo}[2]{}

\newcommand*{\authorsep}{}%
\newcommand*{\lastauthor}{}%
\newcommand*{\prelastauthor}{}%

\newcommand{\displayauthors}{%
  \renewcommand*{\authorsep}{}%
  \renewcommand*{\lastauthor}{}%
  \renewcommand*{\prelastauthor}{}%
  \ifnum\value{NumberOfSuperscripts}=0% If no superscripts are specified, print authors without superscripts.
	\renewcommand*{\do}[1]{%
      \authorsep%
    \lastauthor%
    \renewcommand{\lastauthor}{%
      \renewcommand{\authorsep}{, 
      }%
      \renewcommand{\prelastauthor}{\prelastauthorsep\ \GetTranslation{lastauthorseparator}\ }% Terminated commands with \ to preserve following space
      ##1%
    }%
      }%
    \dolistloop{\listauthors}%
  \else
    \renewcommand*{\dotwo}[2]{%
    \authorsep%
    \lastauthor%
    \renewcommand{\lastauthor}{%
      \renewcommand{\authorsep}{, 
      }%
      \renewcommand{\prelastauthor}{\prelastauthorsep\ \GetTranslation{lastauthorseparator}\ }% Terminated commands with \ to preserve following space
      ##1\textsuperscript{##2}%
    }%
  }%
  \looptwo\listauthors\listsuperscripts%
  \fi
  \prelastauthor \lastauthor%
}

\newcommand{\displayaffiliations}{%
  \ifnum\value{NumberOfSuperscripts}=0% If no superscripts are specified, print affiliations without superscripts.
  \renewcommand*{\do}[1]{%
      ##1\\%
  }%
  \else
  \newcounter{AffiliationNumber}
  \renewcommand*{\do}[1]{%
      \stepcounter{AffiliationNumber}
      \textsuperscript{\arabic{AffiliationNumber}}##1\\%
  }%
  \fi
  \dolistloop{\listaffiliations}%
}

\makeatletter

\def\check@author{%
 \@ifundefined{def@multipleauthors}{%
  \ClassWarningNoLine{apa7}{Author not defined}\authorsnames{Author}}{}
 \@ifundefined{@title}{%
  \ClassWarningNoLine{apa7}{Title not defined}\def\@title{Title}}{}
 \@ifundefined{def@multipleaffils}{%
  \ClassWarningNoLine{apa7}{Affiliation not defined}\authorsaffiliations{Affiliation}}{}
   \@ifundefined{def@stu}{% man mode
   	\@ifundefined{@shorttitle}{%
 	 \ClassWarningNoLine{apa7}{Short title not defined}\def\@shorttitle{INSERT SHORTTITLE COMMAND IN PREAMBLE}}{}
 	\@ifundefined{@abstract}{%
  	\ClassWarningNoLine{apa7}{Abstract not defined}}{}
 	\@ifundefined{@keywords}{%
 	 \ClassInfo{apa7}{Keywords not defined}}{}
	 }{ % stu mode
	 \@ifundefined{@course}{%
 	 \ClassWarningNoLine{apa7}{Course title not defined}\def\@course{Course Title}}{}
 	\@ifundefined{@professor}{%
  	\ClassWarningNoLine{apa7}{Professor not defined}\def\@professor{Professor Name}}{}
 	\@ifundefined{@duedate}{%
 	 \ClassWarningNoLine{apa7}{Due date not defined}\def\@duedate{Due Date}}{}
	 \@ifundefined{@shorttitle}{%
         \def\@shorttitle{}}{}%
	 }
	 \@ifundefined{def@multipleauthors}{}{%
	   \ifnum\value{NumberOfSuperscripts}=0%
	   	% do nothing. If no superscripts are specified, print authors and affiliations without superscripts.
	   \else \ifnum\value{NumberOfAuthors}=\value{NumberOfSuperscripts}%
	         % do nothing
	         \else 
	       		\ClassError{apa7}{Number of Authors does not match number of Superscripts}{Please provide an equal number of authors and superscripts.}
		 \fi
	   \fi
          }%
}

%==== Automatic figure size and orientation determination ====

\newsavebox\gr@box
\newlength\gr@boxwidth
\newlength\gr@boxheight

\newcommand{\fitfigure}[2][0.5]{%
%\ifapamodeman
%{% man
%\sbox\gr@box{\includegraphics[width=\linewidth]{#2}}
%\settowidth{\gr@boxwidth}{\usebox\gr@box}
%\settoheight{\gr@boxheight}{\usebox\gr@box}
%\ifdim\gr@boxheight>\gr@boxwidth{% display upright (h)
% \ifdim\gr@boxheight>\textheight%
%  \centerline{\vspace{\fill}\includegraphics[height=\textheight]{#2}\vspace{\fill}}%
% \else%
%  \usebox\gr@box%
% \fi}
%\else{% display rotated
% \sbox\gr@box{\includegraphics[angle=90,width=\linewidth]{#2}}
% \settoheight{\gr@boxheight}{\usebox\gr@box}
% \ifdim\gr@boxheight>\textheight%
%  \centerline{\vspace{\fill}\includegraphics[angle=90,height=\textheight]{#2}\vspace{\fill}}%
% \else%
%  \usebox\gr@box%
% \fi}%
%\fi
%}{% jou, doc
\sbox\gr@box{\includegraphics[width=\linewidth]{#2}}
\settoheight{\gr@boxheight}{\usebox\gr@box}
\ifdim\gr@boxheight>\textheight%
 \centerline{\includegraphics[height=#1\textheight]{#2}}% need to leave space for caption
\else%
 \usebox\gr@box%
\fi
%}
}

\newcommand{\fitbitmap}[2][0.5]{ % like fitfigure but no scaling in man mode for best quality
%\ifapamodeman
%{% man
%\sbox\gr@box{\includegraphics{#2}}
%\settowidth{\gr@boxwidth}{\usebox\gr@box}
%\settoheight{\gr@boxheight}{\usebox\gr@box}
%\ifdim\gr@boxwidth>\textwidth% display rotated
%  \centerline{\vspace{\fill}\includegraphics[angle=90]{#2}\vspace{\fill}}%
%  \ifdim\gr@boxheight>\textwidth\ClassWarning{apa7}{Figure #2 too high!}\fi% display rotated
%  \ifdim\gr@boxwidth>\textheight\ClassWarning{apa7}{Figure #2 too wide!}\fi% display rotated
%\else% display upright
%  \centerline{\vspace{\fill}\includegraphics{#2}\vspace{\fill}}%
%  \ifdim\gr@boxheight>\textheight\ClassWarning{apa7}{Figure #2 too high!}\fi% display rotated
%\fi}
%{% jou, doc
\sbox\gr@box{\includegraphics[width=\linewidth]{#2}}
\settoheight{\gr@boxheight}{\usebox\gr@box}
\ifdim\gr@boxheight>\textheight%
 \centerline{\includegraphics[height=#1\textheight]{#2}}% need to leave space for caption
\else%
 \centerline{\usebox\gr@box}%
\fi
%}
}

%==== Miscellaneous definitions

\let\normal@BBAB\BBAB                        % -- thp 2005/07/23
\let\table@BBAB\BBAA                         % -- thp 2005/07/23

\setlength{\doublerulesep}{\arrayrulewidth}
\newcommand\thickline{\hline\hline}
\renewcommand\footnoterule{%
  \kern-3\p@
%  \hrule.5\@width.4\columnwidth
  \hrule height0.125pt width.5in
  \kern2.6\p@}

\let\apaSeventabular\tabular
\let\apaSeven@doc@tabular\tabular
%\let\old@tabular\tabular
%\def\apaSeventabular#1{\def\@halignto{to\linewidth}%
% \let\BBAB\table@BBAB%                         -- thp 2005/07/23
% \@tabular{@{\extracolsep{\fill}}#1}}
%\def\apaSeven@doc@tabular{%% thp 2006/01/02
% \let\BBAB\table@BBAB%% Erik Meijer noticed that in doc mode tables weren't centered
% \hfill\old@tabular} %% and that the in-table citation form was incorrect
%%
\let\apaSeven@doc@endtabular\endtabular
%\let\orig@endtabular\endtabular              % -- thp 2005/07/23
%\def\endtabular{%
% \let\BBAB\normal@BBAB%                        -- thp 2005/07/23
% \orig@endtabular\ifapamodedoc{\hfill}{}}% 2006/01/02 \hfill for centering in doc mode
%\def\apaSeven@doc@endtabular{%% thp 2006/01/03
% \let\BBAB\normal@BBAB% %% separate redefinition of tabular/endtabular for doc mode
% \orig@endtabular\hfill}%% to ensure centering, especially in case array.sty is loaded

% Define footnote mark commands for tables that won't mess up vertical alignment
\def\@tab@fn#1{\ensuremath{^{\mbox{{\scriptsize #1}}}}}
\def\tabfnm#1{\rlap{\@tab@fn{#1}}}
%\def\tabfnt#1#2{\par\raggedright\@tab@fn{#1}#2}%BDB
\def\tabfnt#1#2{\raggedright\@tab@fn{#1}#2}

% Provide a command to display vector symbols appropriately in each mode, thp 10/00
% Use improved definition of apaSevenvector provided by Erik Meijer -- 2005/12/28
\def\apaSevenvector#1{{\ensuremath
  \uprightlowercasegreek
  \ifapamodeman
  {\apaSevensmash{\mathop{\kern\z@\mathrm{#1}}\limits_{\scriptscriptstyle\sim}}}%
  {\if@bm@loaded\bm{\mathrm{#1}}\else\mathbf{#1}\fi% in case bm is not available
  }%
}}
% and here are the necessary definitions & packages, as instructed by Erik Meijer
\newcommand{\apaSevensmash}{%
  \def\finsm@sh{\dp\z@\z@ \box\z@}%
  \expandafter\mathpalette\expandafter\mathsm@sh
}%
\newif\if@bm@loaded\@bm@loadedfalse
\IfFileExists{bm.sty}{\RequirePackage{bm}\@bm@loadedtrue}{}% if not, apaSevenvector will fail
\newcommand{\uprightlowercasegreek}{%
  \@ifundefined{alphaup}{}{%
    \def\alpha     {\alphaup     }%
    \def\beta      {\betaup      }%
    \def\gamma     {\gammaup     }%
    \def\delta     {\deltaup     }%
    \def\epsilon   {\epsilonup   }%
    \def\varepsilon{\varepsilonup}%
    \def\zeta      {\zetaup      }%
    \def\eta       {\etaup       }%
    \def\theta     {\thetaup     }%
    \def\vartheta  {\varthetaup  }%
    \def\iota      {\iotaup      }%
    \def\kappa     {\kappaup     }%
    \def\lambda    {\lambdaup    }%
    \def\mu        {\muup        }%
    \def\nu        {\nuup        }%
    \def\xi        {\xiup        }%
    \def\pi        {\piup        }%
    \def\varpi     {\varpiup     }%
    \def\rho       {\rhoup       }%
    \def\varrho    {\varrhoup    }%
    \def\sigma     {\sigmaup     }%
    \def\varsigma  {\varsigmaup  }%
    \def\tau       {\tauup       }%
    \def\upsilon   {\upsilonup   }%
    \def\phi       {\phiup       }%
    \def\varphi    {\varphiup    }%
    \def\chi       {\chiup       }%
    \def\psi       {\psiup       }%
    \def\omega     {\omegaup     }%
  }%
}
\let\apaSevenmatrix\apaSevenvector

%==== Appendix macros added by Michael Erickson 2000/07/11, revised by thp 2000/07/18

\newcounter{appendix}\setcounter{appendix}{0}
\renewcommand{\theappendix}{\@Alph\c@appendix}
\def\apaSevenappfig{%
 \renewcommand\thefigure{\theappendix\@arabic\c@figure}%
 \ifapamodeman{\renewcommand\thepostfig{\theappendix\arabic{postfig}}}{}}
\def\apaSevenapptab{%
 \renewcommand\thetable{\theappendix\@arabic\c@table}%
 \ifapamodeman{\renewcommand\theposttbl{\theappendix\arabic{posttbl}}}{}}
\def\apaSevenappeq{%
 \renewcommand\theequation{\theappendix\@arabic\c@equation}}
\newif\ifoneappendix
\oneappendixtrue % one appendix by default
\newif\ifappendix
\appendixfalse

\def\appendix{%
%
  \ifapamodeman{\processdelayedfloats}{}%  BDB -- output all tables and figures prior to the appendix
%
  \appendixtrue
  \apaSevenappfig
  \apaSevenapptab
  \apaSevenappeq
%  \ifapamodejou{}{\clearpage} %% commented out -- thp 2005/07/23
  \let\old@apaSeven@section=\section % This will not work right with five levels in appendix.
                                 % Should go into \section, not \leveltwo but would also require
                                 % changes to section* and section[ (see \def\section above)
                                 % Who uses five level heading appendices anyway?
%  \newlength{\app@t@width}
%  \setlength{\app@t@width}{\columnwidth}
%  \addtolength{\app@t@width}{-8em}%% leveltwo is defined with leftskip=rightskip=4em plus 1fill
  \long\def\section##1{%
                   \makeatletter%
                     \def\@currentlabelname{##1}%
                   \makeatother%
                   \ifapamodeman{%
                    \clearpage
                    \setcounter{postfig}{0}
                    \setcounter{posttbl}{0}
%                    \efloat@condopen{fff}
%                    \efloat@iwrite{fff}{\string\addtocounter{appendix}{1}}
%                    \efloat@iwrite{fff}{\string\setcounter{figure}{0}}
%                    \efloat@condopen{ttt}
%                    \efloat@iwrite{ttt}{\string\addtocounter{appendix}{1}}
%                    \efloat@iwrite{ttt}{\string\setcounter{table}{0}}%
                   }{%
%                    \setcounter{figure}{0}%
%                    \setcounter{table}{0}%
%                    \vskip2.5ex%
%                    %\small% the appendix should be set in smaller type in jou, including headers!
                   }%
                    \setcounter{figure}{0}%
                    \setcounter{table}{0}%
                    \setcounter{equation}{0}%
                    \vskip2.5ex%
%                  \addtocounter{appendix}{1}%
                   \refstepcounter{appendix}% 2002/07/20 this takes care of references too
                   \ifnum\c@appendix>1\immediate\write\@auxout{\global\string\oneappendixfalse}\fi%
%                   \old@apaSeven@section{%
%                      \old@tabular[t]{@{}p{\app@t@width}@{}}% tabular may have been redefined
%                         \centering{%
%                          \appendixname\ifoneappendix\else~\theappendix\fi\\
%                          ##1%
%                      }\endtabular
%                   }%
%
                      \centerline{\normalfont\normalsize\textbf\appendixname\ifoneappendix\else~\textbf\theappendix\fi}%
                      \centerline{\normalfont\normalsize\bfseries##1}\par%
                      \setlength{\parindent}{0.5in}
                      \makeatletter%
                        \@afterindentfalse%
                        \@afterheading%
                      \makeatother%
%
                  }%
%
} % end of appendix definition

%==== End of general apa.cls macro definitions ====


%%%%%%%%%%%%%%%%%%%%%%%%%
%%                     %%
%%  MANUSCRIPT FORMAT  %%
%%           AND STUDENT        %%
%%                     %%
%%%%%%%%%%%%%%%%%%%%%%%%%

\@ifundefined{def@man}{}{%

% \def\@@spacing{1.75}
\def\@@spacing{1.655}

\newcommand{\@doublespacing}{\linespread{1.655}}
\@doublespacing

% allow a little more space between captions and floats in man mode
\captionsetup[table]{skip=10pt}
\captionsetup[figure]{skip=10pt}

\def\rightheader#1{\def\r@headr{\protect\MakeUppercase{#1}}}
\def\leftheader#1{\def\r@headl{#1}}
%\def\shorttitle#1{\def\s@title{#1}}%\markright{\hfill #1 \thepage} %% no rm - thp 020227

%%%%%%%%%%%%%%%%%%%%%%%%%%%%%%%\newcommand{\shorttitle}[1]{\def\@shorttitle{#1}}

\raggedright

\RequirePackage{fancyhdr}
\setlength{\headheight}{15.2pt}
\fancyhf{}
\renewcommand{\headrulewidth}{0pt}

\fancypagestyle{titlepage}{%
    \lhead{\MakeUppercase{\@shorttitle}}%
    \rhead{\thepage}%
}
\fancypagestyle{otherpage}{%
    \lhead{\MakeUppercase{\@shorttitle}}%
    \rhead{\thepage}%
}
\pagestyle{otherpage}

%========== Alterations to endnotes ===================

%\RequirePackage{endnotes}
%
%% added by BDB to make section heading non-bolface
%\def\enoteheading{\section*{\normalfont\normalsize\notesname
%  \@mkboth{\MakeUppercase{\notesname}}{\MakeUppercase{\notesname}}}%
%  \mbox{}\par\vskip-\baselineskip}
%
%
%\def\@makeenmark{\hbox{$^{\mbox{\small\@ifundefined{hrm}{}{\sf}\@theenmark}}$}}
%\def\enoteformat{\rightskip\z@ \leftskip\z@
%     \raggedright\fussy\hyphenpenalty 10000 \parindent=0.4in
%     \leavevmode\llap{\hbox{$^{\mbox{\small\@ifundefined{hrm}{}{\sf}\@theenmark}}$}}}
%\AtEndDocument{%
% \ifappendix% set section back to no numbering after appendix
%  \appendixfalse\let\section=\old@apaSeven@section%
% \fi
%% \@ifundefined{@acks}{}{\newpage\section{\acksname}\@acks}
% \@ifundefined{hrm}{\def\enotesize{\normalsize}}
%                   {\def\enotesize{\normalsize\@ifundefined{hrm}{}{\sf}}}
% \@ifundefined{NoEndnotes}{\if@enotesopen\newpage\theendnotes\fi}{}
% \let\footnote=\oldfootnote
%}

%========= Alterations to endfloat ====================

% prepare to bypass endfloat for any tables or figures in appendices
\let\@noendfloattab\table% BDB
\let\@noendfloatfig\figure% BDB

% \@ifundefined{float@mark}{% Figure and table in-text markers are no longer required
%  \RequirePackage[notablist,figlist,notabhead,nofighead,tablesfirst,nomarkers]{endfloat}[1995/10/11]}{%
%  \RequirePackage[notablist,figlist,notabhead,nofighead,tablesfirst]{endfloat}[1995/10/11]}

\RequirePackage[notablist,figlist,notabhead,nofighead,tablesfirst,nomarkers]{endfloat}[1995/10/11]

%\renewcommand{\figureplace}{%
%   \par\begin{center}
%   \parbox{\textwidth}
%   {\begin{center}
%   \rule{2.1in}{0.005in}\\
%   Insert \figurename~\thepostfig\ about here\\
%   \rule{2.1in}{0.005in}\\
%   \end{center}}
%   \end{center}\par}
%\renewcommand{\tableplace}{%
%   \par\begin{center}
%   \parbox{\textwidth}
%   {\begin{center}
%   \rule{2.1in}{0.005in}\\
%   Insert \tablename~\theposttbl\ about here\\
%   \rule{2.1in}{0.005in}\\
%   \end{center}}
%   \end{center}\par}

\def\@gobbleuntilnext[#1]{}
\let\eatarg\@gobbleuntilnext
\let\ifnextchar\@ifnextchar

\@ifundefined{def@floatsintext}{%  Guillaume Jourjon 19/10/10
  \def\figure{%
      \ifappendix
          \vspace*{\intextsep}
          \def\fps@figure{!hbt}%
          \@noendfloatfig
      \else
           \efloat@condopen{fff}
           \@nameuse{efloat@separator@fff}%
           \efloat@iwrite{fff}{\string\begin{figure*}[hbt]}%
           \global\def\@figure@written{\relax}% Set a flag that there is at least one figure -- thp 20010705
      %% should be a global declaration to be "visible" at end document
           \ifnextchar[{\gobbleuntilnext[}{}
           \efloat@iwrite{fff}{\string\ifnextchar[{\string\eatarg}{}}
%          \if@domarkers%
%             \addtocounter{postfig}{1}% % bj
%             \figureplace%              % bj
%          \fi%
          \def\@currenvir{efloat@float}%
          \begingroup%
          \let\do\ef@makeinnocent \dospecials%
          \ef@makeinnocent\^^L% and whatever other special cases
          \endlinechar`\^^M \catcode`\^^M=12 \ef@xfigure%
      \fi%
  }%

  \def\table{%
      \ifappendix
          \vspace*{\intextsep}
          \def\fps@table{!hbt}
          \@noendfloattab
      \else
          \efloat@condopen{ttt}
          \@nameuse{efloat@separator@ttt}%
          \efloat@iwrite{ttt}{\string\begin{table*}[hbt]}%
          \ifnextchar[{\gobbleuntilnext[}{}
          \@ifundefined{hrm}{}{%
          \efloat@iwrite{ttt}{\string\sf}}%
          \efloat@iwrite{ttt}{\string\ifnextchar[{\string\eatarg}{}} % bj
%          \if@domarkers
%             \addtocounter{posttbl}{1} % bj
%             \tableplace               % bj
%          \fi
          \def\@currenvir{efloat@float}%
          \begingroup
          \let\do\ef@makeinnocent \dospecials
          \ef@makeinnocent\^^L% and whatever other special cases
          \endlinechar`\^^M \catcode`\^^M=12 \ef@xtable%
      \fi
  }%

% for rotated tables, per Ulrike Fischer:  http://tex.stackexchange.com/a/79843
  \AtEndPreamble{%
    \@ifpackageloaded{rotating}{%
      \DeclareDelayedFloatFlavor{sidewaystable}{table}
      \DeclareDelayedFloatFlavor{sidewaysfigure}{figure}
      }{}%
  }%

}{%
  \RequirePackage{float}
  \floatplacement{figure}{htb}
  \def\figure{\@float{figure}}
  \def\endfigure{\end@float}

  \floatplacement{table}{htb}
  \def\table{\@float{table}}
  \def\endtable{\end@float}
}

%%%%%%%%%%%%%%\long\def\@@contentsline#1#2#3{ #2 }
%%%%%%%%%%%%%%\long\def\numberline#1#2{\noindent{\em\figurename\ #1.\/} #2\vspace{0.5\baselineskip}\par}
%%%%%%%%%%%%%%\long\def\@@caption{\refstepcounter\@captype \@dblarg{\@@@caption\@captype}}
%%%%%%%%%%%%%%\long\def\@@@caption#1[#2]#3{\addcontentsline{\csname
%%%%%%%%%%%%%%  ext@#1\endcsname}{#1}{\protect\numberline{\csname
%%%%%%%%%%%%%%  the#1\endcsname}{\ignorespaces #2}}}
%\def\listoffigures{\section*{\figurecaptionsname}\@starttoc{lof}}% changed from 'Figure Captions'

%\long\def\figurecaptionsformat{
%  \let\contentsline=\@@contentsline
%  \let\numberline=\@@numberline
%  \def\baselinestretch{\@@spacing}
%}

\def\processfigures{%
 \expandafter\ifnum \csname @ef@fffopen\endcsname>0
  \immediate\closeout\efloat@postfff \ef@setct{fff}{0}
  \clearpage
  \if@figlist
   \@ifundefined{@figure@written}{}{%
%    \figurecaptionsformat
    {\normalsize\@ifundefined{hrm}{}{\sf}%
%     \let\BBAB\table@BBAB% -- thp 2005/07/23
%     \listoffigures%
%     \let\BBAB\normal@BBAB% -- thp 2005/07/23
    }
%    \clearpage
   }
  \fi
%  \if@fighead
%     \section*{\figuresection}
%     \suppressfloats[t]
%  \fi
  \@ifundefined{fig@num}{%
%  \markboth{}{}
   \pagestyle{otherpage}
%  \pagestyle{empty}%
  }{%
  \setcounter{page}{1}
%  \markright{\rm\hfill \s@title, Figure}
   \def\@oddhead{\rightmark}                                    % changed by Michael Erickson
   \markright{\hfill \s@title, \figurename\ \protect\thefigure} % to include appendix numbering
                                                                % remove rm - thp 020227
  }
%  \let\caption=\@@caption
  \def\baselinestretch{\@@spacing}\normalsize\@ifundefined{hrm}{}{\sf}
  \processfigures@hook \@input{\jobname.fff}
 \fi}

\def\processtables{%
  \expandafter\ifnum \csname @ef@tttopen\endcsname>0
  \immediate\closeout\efloat@postttt \ef@setct{ttt}{0}
  \clearpage
  \if@tabhead
      \section*{\tablesection}
      \suppressfloats[t]
  \fi
  \def\baselinestretch{\@@spacing}
  \processtables@hook \@ifundefined{hrm}{}{\sf}%
  \tiny\normalsize%
% We need to deal with citation alterations here because otherwise we either
% lose the definition through the .ttt diversion or we miss the table*
  \let\BBAB\table@BBAB%  -- thp 2005/07/23
  \@input{\jobname.ttt}%
  \let\BBAB\normal@BBAB% -- thp 2005/07/23
 \fi}

%========== Alterations to captions (BDB) ===================

\captionsetup{justification=raggedright}


%========= Maketitle, margins, lengths, etc. ==========

\def\maketitle{
\@ifundefined{hrm}{}{\hrm}
 \check@author
% \@ifundefined{s@title}{\ClassInfo{apa7}{Using title for short title}
%                        \def\s@title{\@title}
%                        \markright{\hfill\@title}}%\protect\thepage}}  %% removed rm
%                       {\markright{\hfill\s@title}}%\protect\thepage}} %% thp 020227
% \@ifundefined{r@headr}{\ClassInfo{apa7}{Using short title for running head}
%                        \def\r@headr{\protect\MakeUppercase{\s@title}}}{}


  \begin{center}
  
%  \makebox[\linewidth][l]{\rheadname\hspace{0.1in}\MakeUppercase{\r@headr}}\\ %thp090298

  \vspace*{4\baselineskip}
  \textbf\@title%
  \ifapamodeman{%
    \vspace{\baselineskip}\\
  }{%
    \vspace{\baselineskip}\\
  }
\@ifundefined{def@stu}{ %Professional manuscript
  \@ifundefined{apaSeven@maskauthoridentity}{%  BDB
      
      \displayauthors \\
      \displayaffiliations

      \@ifundefined{@note}
       {\vspace*{\baselineskip} }
       {\@note}

  }{%  mask author identity -- show nothing in the author or author note space
  }

  \end{center}

  \@ifundefined{apaSeven@maskauthoridentity}{
      \@ifundefined{@acks}
       {}
       {%
         \vfill%
        \section*{\normalfont\normalsize\textbf\noexpand\GetTranslation{acksname}}
         \raggedright
         \setlength{\parindent}{0.5in}
         \indent\@acks\par%
       }
  }{%  mask author identity -- show nothing in the author or author note space
  }}{% Student Manuscript
    \@ifundefined{apaSeven@maskauthoridentity}{%  BDB
      \displayauthors \\
      
      \displayaffiliations
    }{%  mask author identity -- show nothing in the author or affiliations space
    }
      \@course \\ 
      \@professor \\ 
      \@duedate \\ 
      
    \@ifundefined{apaSeven@maskauthoridentity}{%  BDB
      \@ifundefined{@note}
       {\vspace*{\baselineskip} }
       {\@note}
    }{%  mask author identity -- show nothing in the note space
    }
       
        \end{center}
       
       \@ifundefined{@acks}
       {}
       {%
         \@ifundefined{apaSeven@maskauthoridentity}{%  BDB
           \vfill%
           \section*{\normalfont\normalsize\textbf\noexpand\GetTranslation{acksname}}
           \raggedright
           \setlength{\parindent}{0.5in}
           \indent\@acks\par%
         }{%  mask author identity -- show nothing in the author note space
         }
       }
  
  }
  \newpage
  %BDB\hyphenpenalty 10000
  \fussy
  \@ifundefined{@abstract}{}{%
    \section*{\normalfont\normalsize\textbf\abstractname}% BDB
    \noindent\@abstract\par% BDB
    \@ifundefined{@keywords}{}{%
      \setlength{\parindent}{0.5in}% BDB
      \indent\textit{\GetTranslation{keywordname}:} \@keywords%
    }%
    \newpage
  }

  \@ifundefined{def@donotrepeattitle}{
    \section*{\protect\normalfont\textbf{\@title}}
  }{}%
  \raggedright%
  \setlength{\parindent}{0.5in}%
%  \indent%
}

\thispagestyle{titlepage}

%\ThreeLevelHeading
% \def\baselinestretch{\@@spacing}

% footnote issues: separate multiple footnotes on the same page:
\setlength{\footnotesep}{16pt}

% make footnotes ragged-right (credit to Mako Hill):
\renewcommand\@makefntext[1]{\raggedright\textsuperscript{\@thefnmark}~#1}

%\setlength{\topmargin}{0in}
%\setlength{\oddsidemargin}{0.25in}
%\setlength{\evensidemargin}{0.25in}
%\setlength{\textwidth}{6in}
%\setlength{\textheight}{8.5in}
%\setcounter{secnumdepth}{0}
%
%\def\ps@myheadings{
% \let\@mkboth\@gobbletwo
% \def\@oddhead{\hbox{}\rightmark\hfil\hspace{0.1in}\thepage} %% removed rm - thp 020227
% \def\@oddfoot{}
% \def\@evenhead{}
% \def\@evenfoot{}
% \def\sectionmark##1{}
% \def\subsectionmark##1{}}

%\pagestyle{myheadings}
%\raggedright
%BDB\hyphenpenalty 10000
\newcommand{\footmark}[1]{${}^{\mbox{\normalsize #1}}$}

%\long\def\@makecaption#1#2{%
% \let\BBAB\table@BBAB % thp 2005/07/23
% \vskip10pt%
% \setbox\@tempboxa\hbox{\parbox{\columnwidth}{#1\cap@line{\cap@style #2}}}%
% \ifdim\wd\@tempboxa>\hsize\parbox{\columnwidth}{#1\cap@line{\cap@style #2}}\par%
% \else\hbox to\hsize{\box\@tempboxa\hfil}\par%
% \fi%
% \let\BBAB\normal@BBAB % thp 2005/07/23
% \vskip\baselineskip}
%
%\long\def\@caption#1[#2]#3{\par\addcontentsline{\csname
%  ext@#1\endcsname}{#1}{\protect\numberline{\csname
%  the#1\endcsname}{\ignorespaces #2}}\begingroup
%    \@parboxrestore
%    \normalsize\@ifundefined{hrm}{}{\hrm}
%    \@makecaption{\csname fnum@#1\endcsname}{\ignorespaces #3 } \par
%  \endgroup}
%
%\def\caption{\refstepcounter\@captype \@dblarg{\@caption\@captype}}
%
%\def\centeredcaption#1{\caption}
%
%\def\fnum@figure{\def\cap@style{}\def\cap@line{}{\em\figurename\ {\thefigure}}. }
%% added second set of braces around \em  to get citations in man mode -- tp 17/7/2000
%% then removed them again because they were cancelling application of em to the caption
%\long\def\fnum@table{\def\cap@style{\em}\def\cap@line{\\ }\tablename\ \thetable }

\setcounter{topnumber}{1}
\def\topfraction{.7}
\setcounter{bottomnumber}{1}
\def\bottomfraction{.6}
\setcounter{totalnumber}{1}
\def\textfraction{0}
\def\floatpagefraction{.7}
\setcounter{dbltopnumber}{1}
\def\dbltopfraction{.7}
\def\dblfloatpagefraction{.7}
\def\dbltextfloatsep{\textfloatsep}

\fussy
%\setlength{\parindent}{0.4in}
%\let\oldfootnote=\footnote % 3/10/00
%\let\footnote=\endnote
%\def\footnotesize{\normalsize\rm}

\@ifundefined{def@nosf}{%
\def\helvetica{%
\ClassWarning{apa7}{ignored \string\helvetica\space (use helv option)}
}}{\def\helvetica{\relax}}

%==== Appendix definitions added by Michael Erickson 2000/07/11

%\efloat@condopen{fff}
%\efloat@iwrite{fff}{\string\apaSevenappfig}
%\efloat@iwrite{fff}{\string\setcounter{appendix}{0}}
%\efloat@condopen{ttt}
%\efloat@iwrite{ttt}{\string\apaSevenapptab}
%\efloat@iwrite{ttt}{\string\setcounter{appendix}{0}}

}% end of man mode (manuscript format)


%%%%%%%%%%%%%%%%%%%%
%%                %%
%% JOURNAL FORMAT %%
%%                %%
%%%%%%%%%%%%%%%%%%%%

\@ifundefined{def@jou}{}{%

\RequirePackage{fancyhdr}

\renewcommand{\addORCIDlink}[2]{\href{https://orcid.org/#2}{{\mbox{\scalerel*{
\begin{tikzpicture}[yscale=-1,transform shape]
\pic{orcidlogo};
\end{tikzpicture}
}{|}}} #1}}

%========== balance last-page columns =================

\IfFileExists{flushend.sty}{\RequirePackage{flushend}}{}

%========== ftnright and alterations ==================

\IfFileExists{ftnright.sty}{
 \let\savefootnoterule\footnoterule
 \let\save@makefntext\@makefntext
 \RequirePackage{ftnright}
 \let\footnoterule\savefootnoterule
 \let\@makefntext\save@makefntext
}{}

%======================================================

\def\rightheader#1{\def\r@headr{\protect\MakeUppercase{\protect\scriptsize #1}}}
\def\leftheader#1{\def\r@headl{\protect\MakeUppercase{\protect\scriptsize #1}}}
\def\r@headr{\protect\MakeUppercase{\protect\scriptsize\@shorttitle}}% BDB
%%%%%%%%%%%%%%%%%%%%%%%\def\shorttitle#1{\def\r@headr{\protect\MakeUppercase{\protect\scriptsize #1}}}% BDB

\def\maketitle{
 \check@author
 \@ifundefined{r@headr}{\def\r@headr{\protect\MakeUppercase{\protect\scriptsize\@title}}}{}
 \@ifundefined{r@headl}{\def\r@headl{ADD \\LEFTHEADER COMMAND}}{}

\twocolumn[  % anything appearing within the brackets is set in one-column mode
  \vspace{0.03in}
  \begin{center}
% title
  {\LARGE \@title}\\
  \vspace{-0.05in}

  \@ifundefined{apaSeven@maskauthoridentity}{%  BDB

	\parbox[t]{\textwidth}{\begin{center}{\large \displayauthors\vspace{0in}}%

	{\displayaffiliations\vspace{0.05in}}\end{center}}
	
      \@ifundefined{@note}
       {\vspace{0.07in}}
       {\vspace{0.07in}\\ {\large\@note\vspace{0.07in}}}

  }{%  mask author identity -- show nothing in the author or author space
    \vspace{0.32in}
  }


  \@ifundefined{@abstract}
  {\par }
  {\par \parbox{4.6875in}
   {\small \noindent \@abstract
     \@ifundefined{@keywords}{}{%
      \par\vspace{0.12in}\raggedright\textit{\GetTranslation{keywordname}:} \@keywords%
     }%
   }
   \vspace{0.24in}%
  }
  \end{center}
 ] % end of \twocolumn[]
 
 \renewcommand{\headrulewidth}{0pt}%

 \pagenumbering{arabic}
 \@ifundefined{@journal}{\thispagestyle{empty}}{%
  \@ifundefined{@vvolume}{\def\@vvolume{\strut}}{}%
  \@ifundefined{@copnum}{\def\@copnum{\strut}}{}%
  \@ifundefined{@ccoppy}{\def\@ccoppy{\strut}}{}%
  \fancypagestyle{pageone}{%
    \fancyhf{}%
    \fancyhead[L]{\stiny{\@journal}\vspace{-0.15\baselineskip}\\
                 \stiny{\@vvolume}}%
    \fancyhead[R]{\stiny{\@ccoppy}\vspace{-0.15\baselineskip}\\
                 \stiny{\@copnum}}%
  }
  \thispagestyle{pageone}
 }

  \@ifundefined{apaSeven@maskauthoridentity}{%  BDB
     \@ifundefined{@acks}
      {}
      {\begin{figure}[b]
       \parbox{\columnwidth}{\setlength{\parindent}{0.18in}
       \noindent\makebox[\columnwidth]{\vrule height0.125pt width\columnwidth}\vspace*{0.05in}\par
       {\footnotesize\hspace{-0.04in}\@acks\par}}
       \end{figure}}
  }{%  mask author identity -- show nothing in the author note space
  }

 \@ifundefined{no@tab}{\let\tabular\apaSeventabular}{}
 %\noindent
}

\newcommand\stiny{\@setfontsize\stiny\@vipt\@viipt}

%\ThreeLevelHeading
%\def\baselinestretch{0.923}
%\def\refname{References}
\setlength{\footnotesep}{0.2813in}
\setlength{\topmargin}{-0.275in}
\addtolength{\headheight}{0.02in}
\addtolength{\headsep}{-0.156in}
\setlength{\oddsidemargin}{-0.25in}
\setlength{\evensidemargin}{-0.25in}
\setlength{\textwidth}{6.94in}
\setlength{\textheight}{8.9in}
\setlength{\columnwidth}{8.5cm}
\setlength{\columnsep}{0.25in}
\setlength{\parindent}{0.15625in}
%%\setlength{\parskip}{0in}
\setlength{\textfloatsep}{0.35in}
%\flushbottom

\setcounter{secnumdepth}{0}

 \fancypagestyle{journalbase}{%
  \fancyhf{}%
  \fancyhead[LE]{\thepage}%
  \fancyhead[RO]{\thepage}%
  \fancyhead[CO]{\r@headr}%
  \def\sectionmark##1{}
  \def\subsectionmark##1{}
}

\@ifundefined{apaSeven@maskauthoridentity}{%  BDB
    \fancypagestyle{journal}[journalbase]{%
      \fancyhead[CE]{\r@headl}%
    }
}{%  mask author identity -- show the short title for both the left and right headers
    \fancypagestyle{journal}[journalbase]{%
      \fancyhead[CE]{\r@headr}%
    }
 }
 \pagestyle{journal}

%\long\def\@makecaption#1#2{%
%  \let\BBAB\table@BBAB % thp 2005/07/23
%  \vskip\abovecaptionskip
%  \sbox\@tempboxa{#1 #2}%
%  \ifdim \wd\@tempboxa >\hsize
%    #1 #2\par
%  \else
%    \global \@minipagefalse
%    \hb@xt@\hsize{\hfil\box\@tempboxa\hfil}%
%  \fi
%  \let\BBAB\normal@BBAB % thp 2005/07/23
%  \vskip\belowcaptionskip}
%
%\long\def\@makeccaption#1#2#3{
% \let\BBAB\table@BBAB % thp 2005/07/23
% \vskip\abovecaptionskip
% \setbox\@tempboxa\hbox{#1 #2}
% \ifdim \wd\@tempboxa > \hsize
%    \hbox to\hsize{\hfil\parbox[t]{#3}{#1 #2}\hfil}
%  \else
%    \ifdim \wd\@tempboxa > #3
%      \hbox to\hsize{\hfil\parbox[t]{#3}{#1 #2}\hfil}
%    \else
%      \hbox to\hsize{\hfil\mbox{#1 #2}\hfil}
%    \fi
% \fi
% \let\BBAB\normal@BBAB % thp 2005/07/23
% \vskip\belowcaptionskip}
%
%\long\def\@caption#1[#2]#3{\par\addcontentsline{\csname
%  ext@#1\endcsname}{#1}{\protect\numberline{\csname
%  the#1\endcsname}{\ignorespaces #2}}\begingroup
%    \@parboxrestore
%    \normalsize
%    \csname setcskip@#1\endcsname
%    \@makecaption{\csname fnum@#1\endcsname}{\ignorespaces #3}\par
%  \endgroup}
%
%\long\def\@ccaption#1[#2]#3#4{\def\@ccwidth{#3} \par\addcontentsline{\csname
%  ext@#1\endcsname}{#1}{\protect\numberline{\csname
%  the#1\endcsname}{\ignorespaces #2}}\begingroup
%    \@parboxrestore
%    \normalsize
%    \@makeccaption{\csname fnum@#1\endcsname}{\ignorespaces #4}{\@ccwidth}\par
%  \endgroup}
%
%\def\caption{\refstepcounter\@captype \@dblarg{\@caption\@captype}}
%
%\long\def\centeredcaption{\refstepcounter\@captype \@dblarg{\@ccaption\@captype}}
%
%\def\setcskip@table{\setlength{\abovecaptionskip}{3.2pt}\setlength{\belowcaptionskip}{2pt}}
%\def\setcskip@figure{\setlength{\abovecaptionskip}{-0.04in}\setlength{\belowcaptionskip}{-1.9pt}}
%
%\def\fnum@figure{\small{\em\figurename\ {\thefigure}}.\hspace{0.07in}}
%\long\def\fnum@table{\makebox[\linewidth][l]{\tablename\ \thetable} \\ \em }
%
%\def\@makefnmark{\hbox{\@textsuperscript{\normalfont{\scriptsize\@thefnmark}}}}%
%\long\def\@makefntext#1{\parindent 1em\noindent
%           \hb@xt@1.8em{%
%           \hss\@textsuperscript{\normalfont{\tiny\@thefnmark}\hspace{1.5pt}}}#1}%


\setcounter{topnumber}{2}
\def\topfraction{.85}
\setcounter{bottomnumber}{2}
\def\bottomfraction{.75}
\setcounter{totalnumber}{3}
\def\textfraction{.10}
\def\floatpagefraction{.85}
\setcounter{dbltopnumber}{2}
\def\dbltopfraction{.85}
\def\dblfloatpagefraction{.85}
\def\dbltextfloatsep{0.8\textfloatsep}
\let\footnotesize=\small

\def\helvetica{\relax}

\doublehyphendemerits5000
\hfuzz0pt
\tolerance=9999
\pretolerance=-1
\emergencystretch=25pt
\hbadness=30000
\hyphenpenalty=100


\@ifundefined{def@apacite}{}{% -- Philip Kime 2008/12/03
  % Removed some bibliography redefinitions as per the instructions of Erik Meijer
  \bibleftmargin=1.2em        % left these in because the default is too big
  \bibindent=-\bibleftmargin  % and this is apparently not refefined each time
  \renewcommand{\bibliographytypesize}{\footnotesize}}
\@ifundefined{def@natbib}{}{% -- Philip Kime 2008/12/03
  % Removed some bibliography redefinitions as per the instructions of Erik Meijer
  \bibleftmargin=1.2em        % left these in because the default is too big
  \bibindent=-\bibleftmargin  % and this is apparently not refefined each time
  \renewcommand{\bibliographytypesize}{\footnotesize}}



%\@ifundefined{def@noapacite}{% -- Philip Kime 2008/12/03
%  % Removed some bibliography redefinitions as per the instructions of Erik Meijer
%  \bibleftmargin=1.2em        % left these in because the default is too big
%  \bibindent=-\bibleftmargin  % and this is apparently not refefined each time
%  \renewcommand{\bibliographytypesize}{\footnotesize}}{}

}% end of jou mode (journal format)

%%%%%%%%%%%%%%%%%%%%
%%                %%
%% REGULAR FORMAT %%
%%                %%
%%%%%%%%%%%%%%%%%%%%

\@ifundefined{def@doc}{}{%

%\pagestyle{myheadings}

%%%%%%%%%%%%%%%%%%%%%%%%%%%%%%%%%%%%%%%%%%%%\renewcommand{\shorttitle}[1]{\def\@shorttitle{#1}}

\def\leftheader#1{\def\r@headl{#1}}

\RequirePackage{fancyhdr}
\setlength{\headheight}{15.2pt}
\fancyhf{}
\renewcommand{\headrulewidth}{0pt}

\fancypagestyle{otherpage}{%
% %%%%%%%%%%%%%%%%%%%%%%%%%%%%%%%%%%%%%%%%%%%%%%%%%%%%%%%%%%%%%%%%%%%%%%%%%%%%%%%%%%%%%%%%%%%%%%%%%%%%%%%%%%%%%%%%%%%%%%%
    \lhead{\MakeUppercase{\@shorttitle}}%
    \rhead{\thepage}%
}
\pagestyle{otherpage}

\def\maketitle{
 \global\@topnum\z@  % to prevent tables before the title   %% Erik Meijer, 2006/01/03
 \@ifundefined{@acks}{% if there acknowledgements they make up a "float" on the 1st page
 \global\@botnum\z@}{% to prevent tables on the first page  %% Erik Meijer, 2006/01/03
 \global\@botnum\@ne}% to prevent tables below the footnote -- thp 2006/01/10
 \check@author
% %%%%%%%%%%%%%%%%%%%%%%%%%%%%%%%%%%%%%%%%%%%%%%%%%%%%%%%%%%%%%%%%%%%%%%%%%%%%%%%%%%%%%%%%%%%%%%%%%%%%%%%%%%%%%%%%%%%%%%%
 \@ifundefined{r@headr}{\typeout{Using title for running head}
                        \def\r@headr{\protect\MakeUppercase{\@title}}
                        \markright{\rm \@title \protect\\ \thepage}}{}
 \@ifundefined{r@headl}{\let\r@headl\r@headr}{}
 \@ifundefined{s@title}{\let\s@title\r@headr}{}
  \sloppy
  \setlength{\parindent}{0.5in}
  \begin{center}
   \@ifundefined{@journal}{}{%
    \@ifundefined{@vvolume}{\def\@vvolume{}}{}%
    \@ifundefined{@copnum}{\def\@copnum{}}{}%
    \@ifundefined{@ccoppy}{\def\@ccoppy{}}{}%
    {\scriptsize{\@journal}}\hspace{\fill}{\scriptsize{\@ccoppy}}\vspace{-0.3\baselineskip}\\
    {\scriptsize{\@vvolume}}\hspace{\fill}{\scriptsize{\@copnum}}\vspace{0.1in}\\
   }
  \vspace*{0.3in}

  {\LARGE \@title}\\

  \vspace{0.3in}
  \@ifundefined{apaSeven@maskauthoridentity}{%  BDB

      {\Large \displayauthors} \\

      \displayaffiliations \vspace{0.1in}

  }{%  mask author identity -- show nothing in the author note space
  }

  \@ifundefined{@abstract}{}{%
   {\abstractname}\vspace{0.1in}% BDB

   \parbox{5in}{\@abstract%
     \@ifundefined{@keywords}{}{%
       \par\vspace{0.12in}\textit{\GetTranslation{keywordname}:} \@keywords%
     }%
   }\vspace{0.25in}%
  }
  \end{center}
 \pagenumbering{arabic}
 \thispagestyle{empty}

  \@ifundefined{apaSeven@maskauthoridentity}{%  BDB

     \@ifundefined{@acks}
      {}
      {\begin{figure}[b]
       \parbox{\textwidth}{ \setlength{\parindent}{0.2in}
       \noindent \makebox[\linewidth]{\vrule height0.125pt width\linewidth}

       \vspace*{0.05in}
       {\footnotesize
       \indent \@acks

       }}
       \end{figure}}

  }{%  mask author identity -- show nothing in the author note space
  }

%%%%%%%%%%%%%%%%%%%%%%%%%%%%%%%%%%%%%%%%%%%%%%%%%%%%%%%%%%%%%%%%%%%%%%%%%%%%%%%%%%%%%%%%%%%%%%%%%%%%%%%%%%%%%%%%%%%%%%%%
 \markboth{\hfill \r@headl \hfill}{\hfill \r@headr \hfill}
 \let\tabular\apaSeven@doc@tabular% -- thp 2006/01/02
 \let\endtabular\apaSeven@doc@endtabular% -- thp 2006/01/03
% \noindent
}

%\ThreeLevelHeading
%\def\refname{References}
\setlength{\topmargin}{-0.25in}
\setlength{\oddsidemargin}{0.25in}
\setlength{\evensidemargin}{0.25in}
\setlength{\textwidth}{6in}
\setlength{\textheight}{8.5in}
\setcounter{secnumdepth}{0}
\flushbottom

% fix headwidth for even-numbered pages
% suggestion from http://tex.stackexchange.com/questions/42798/apa7-package-heading-line-exceeds-textwidth
\setlength{\headwidth}{\textwidth}

\setcounter{topnumber}{2}
\def\topfraction{.7}
\setcounter{bottomnumber}{2}
\def\bottomfraction{.6}
\setcounter{totalnumber}{3}
\def\textfraction{.2}
\def\floatpagefraction{.7}
\setcounter{dbltopnumber}{2}
\def\dbltopfraction{.8}
\def\dblfloatpagefraction{.8}
\def\dbltextfloatsep{0.8\textfloatsep}

\def\helvetica{\relax}
%%\def\timesroman{\relax}%    Commented out all \timesroman -- thp 2005/07/23
%\def\centeredcaption#1{\caption}
%
%\long\def\@makecaption#1#2{
% \let\BBAB\table@BBAB% thp 2005/07/23
% \vskip 10pt
% \setbox\@tempboxa\hbox{#1 #2}
% \ifdim \wd\@tempboxa >\hsize #1 #2\par \else \hbox
%to\hsize{\hfil\box\@tempboxa\hfil}
% \fi
% \let\BBAB\normal@BBAB% thp 2005/07/23
% \vskip 4pt}
%
%\long\def\@caption#1[#2]#3{\par\addcontentsline{\csname
%  ext@#1\endcsname}{#1}{\protect\numberline{\csname
%  the#1\endcsname}{\ignorespaces #2}}\begingroup
%    \@parboxrestore
%    \normalsize
%    \@makecaption{\csname fnum@#1\endcsname}{\ignorespaces #3}\par
%  \endgroup}
%
%\def\fnum@figure{\small {\em\figurename\ {\thefigure}}.\hspace{0.07in}}
%\long\def\fnum@table{\small\tablename\ {\thetable}:}


\@ifundefined{def@apacite}{}{% -- Philip Kime 2008/12/03
  % Removed defs as per the instructions of Erik Meijer -- thp 2005/07/23
  \renewcommand{\bibliographytypesize}{\small}}
\@ifundefined{def@natbib}{}{% -- Philip Kime 2008/12/03
  % Removed defs as per the instructions of Erik Meijer -- thp 2005/07/23
  \renewcommand{\bibliographytypesize}{\small}}

%\@ifundefined{def@noapacite}{% -- Philip Kime 2008/12/03
%  % Removed defs as per the instructions of Erik Meijer -- thp 2005/07/23
%  \renewcommand{\bibliographytypesize}{\small}}{}


}% end of doc mode (regular LaTeX format)

%---------- end of mode specific definitions -----------

% =======================================================
%       Some final definitions for all modes
% =======================================================

\let\ignore\@gobble


\@ifundefined{def@apacite}{}{% -- Philip Kime 2008/12/03
  \bibliographystyle{apacite}
  %
  % Thanks to Donald Arsenau for the right way to ignore \bibliographystyle
  %
  \def\bibliographystyle#1{\ClassWarning{apa7}{\string\bibliographystyle\space
      command ignored}}}

\@ifundefined{def@natbib}{}{% -- Philip Kime 2008/12/03
  \bibliographystyle{apacite}
  %
  % Thanks to Donald Arsenau for the right way to ignore \bibliographystyle
  %
  \def\bibliographystyle#1{\ClassWarning{apa7}{\string\bibliographystyle\space
      command ignored}}}

%\@ifundefined{def@noapacite}{% -- Philip Kime 2008/12/03
%  \bibliographystyle{apacite}
%  %
%  % Thanks to Donald Arsenau for the right way to ignore \bibliographystyle
%  %
%  \def\bibliographystyle#1{\ClassWarning{apa7}{\string\bibliographystyle\space
%      command ignored}}}{}




%    \end{macrocode}
%\end{macro}
%
%    \begin{macrocode}
%</class>
%    \end{macrocode}
%
%
% %%%%%%%%%%%  SAMPLE FILES  %%%%%%%%%%%%
%
%\begin{macro}{bibliography.bib}
%    \begin{macrocode}
%<*bibliography>

@ARTICLE{Borst2011a,
  author = {Borst, Gr\'{e}goire and Kosslyn, Stephen M. and Kievit, Rogier A. and  Thompson, William L.},
  title = {A fictitious entry to demonstrate complicated bibliographies},
  journal = {Journal of Nothingness},
  year = {2011},
  volume = {1},
  pages = {2--3}
}

@ARTICLE{Borst2011b,
  author = {Borst, Gr\'{e}goire and Kievit, Rogier A. and Thompson, William L. and
    Kosslyn, Stephen M.},
  title = {Mental rotation is not easily cognitively penetrable},
  journal = {Journal of Cognitive Psychology},
  year = {2011},
  volume = {23},
  pages = {60--75}
}

@ARTICLE{Borst2011c,
  author = {Borst, Gr\'{e}goire and Thompson, William L. and Kosslyn, Stephen M.},
  title = {Understanding the dorsal and ventral systems of the human cerebral
	cortex: {B}eyond dichotomies},
  journal = {American Psychologist},
  year = {2011},
  volume = {66},
  pages = {624--632}
}

@ARTICLE{deWaal2009,
  author = {de Waal, Elda and Grosser, M. M.},
  title = {Safety and security at school: {A} pedagogical perspective},
  journal = {Teaching and Teacher Education},
  year = {2009},
  volume = {25},
  pages = {697--706}
}

@ARTICLE{Franklin2010,
  author = {Franklin, Jr., Robert G. and Adams, Jr., Reginald B.},
  title = {The two sides of beauty: {L}aterality and the duality of facial attractiveness},
  journal = {Brain and Cognition},
  year = {2010},
  volume = {72},
  pages = {300--305}
}

@ARTICLE{Gilbert2004,
  author = {Gilbert, David and McClernon, Joseph and Rabinovich, Norka and Sugai, Chihiro
    and Plath, Louisette and Asgaard, Greg and Zuo, Yantao and	Huggenvik, Jodi and Botros, Nazeih},
  title = {Effects of quitting smoking on {EEG} activation and attention last for more than 31
    days and are more severe with stress, dependence,	{DRD2 A1} allele, and depressive traits},
  journal = {Nicotine \& Tobacco Research},
  year = {2004},
  volume = {6},
  pages = {249--267},
  number = {2},
  doi = {10.1080/14622200410001676305}
}

@INBOOK{Haybron2008,
  author = {Haybron, Daniel M.},
  title = {Philosophy and the science of subjective well-being},
  editor = {Eid, Michael and Larsen, Randy J.},
  booktitle = {The science of subjective well-being},
  year = {2008},
  pages = {17--43},
  publisher = {Guilford Press},
  address = {New York, NY}
}

@ARTICLE{Herbst-Damm2005,
  author = {Herbst-Damm, Kathryn L. and Kulik, James A.},
  title = {Volunteer support, marital status, and the survival times of terminally ill patients},
  journal = {Health Psychology},
  year = {2005},
  volume = {24},
  pages = {225--229},
  number = {2}
}

@ARTICLE{Lassen2006,
  author = {Lassen, Stephen R. and Steele, Michael M. and Sailor, Wayne},
  title = {The relationship of school-wide positive behavior support to academic achievement in
    an urban middle school},
  journal = {Psychology in the Schools},
  year = {2006},
  volume = {43},
  pages = {701--712},
  number = {6}
}

@ARTICLE{JRLevin2000,
  author = {Levin, Joel R. and O'Donnell, Angela M.},
  title = {What to do about educational research's credibility gaps?},
  journal = {Journal of Cognitive Education and Psychology},
  year = {2000},
  volume = {1},
  pages = {201--256}
}

@ARTICLE{MLevin1990,
  author = {Levin, Mary E. and Levin, Joel R.},
  title = {Scientific mnemonomies: {M}ethods for maximizing more than memory},
  journal = {American Educational Research Journal},
  year = {1990},
  volume = {27},
  pages = {301--321}
}

@ARTICLE{MayerInPressA,
  author = {Mayer, Richard E.},
  title = {Fake article to illustrate in-press citations},
  journal = {Journal of Cognitive Psychology},
  year = {in press}
}

@ARTICLE{MayerInPressB,
  author = {Mayer, Richard E.},
  title = {Second fake article to illustrate in-press citations},
  journal = {Journal of Educational Psychology},
  year = {in press}
}

@ARTICLE{Mayer2008a,
  author = {Mayer, Richard E.},
  title = {Applying the science of learning: {E}vidence-based principles for
	the design of multimedia instruction},
  journal = {American Psychologist},
  year = {2008},
  volume = {63},
  pages = {760--769}
}

@ARTICLE{Mayer2008b,
  author = {Mayer, Richard E.},
  title = {Old advice for new researchers},
  journal = {Educational Psychology Review},
  year = {2008},
  volume = {20},
  pages = {19--28}
}

@BOOK{Shotton1989,
  author = {Shotton, Margaret A.},
  title = {Computer addiction? {A} study of computer dependency},
  publisher = {Taylor \& Francis},
  year = {1989},
  location = {London, England}
}

@ARTICLE{vonDavier2011,
  author = {von Davier, Matthias and Xu, Xueli and Carstensen, Claus H.},
  title = {Measuring growth in a longitudinal large-scale assessment with a
    general latent variable model},
  journal = {Psychometrika},
  year = {2011},
  volume = {76},
  pages = {318--336}
}

%</bibliography>
%    \end{macrocode}
%\end{macro}
%
%\begin{macro}{shortsample.tex}
%    \begin{macrocode}
%<*shortsample>
\documentclass[jou]{apa7}

\usepackage[hungarian,british]{babel}

\usepackage{csquotes}
\usepackage[style=apa,backend=biber]{biblatex}
\addbibresource{bibliography.bib}

\title{Sample APA-Style Document Using the \textsf{apa7} Package}

\authorsnames{Daniel A. Weiss}
\authorsaffiliations{Department of Psychology, A University Somewhere}

\leftheader{Weiss}

\abstract{This demonstration paper uses the \textsf{apa7} \LaTeX\
  class to format the document in compliance with the 7th Edition of
  the American Psychological Assocation's \textit{Publication Manual.}
  The references are managed using \textsf{biblatex}.}

\keywords{APA style, demonstration}

\authornote{\addORCIDlink{Daniel A. Weiss}{0000-0000-0000-0000}, Department of Psychology, A University Somewhere

Changes of Affiliations and other info here to describe what needs to be said}

\begin{document}
\maketitle
We begin with \textcite{Shotton1989}.  We can also cite this work in
parenthesis, like this: \parencite{Lassen2006}.

Note the use of five heading levels throughout this demonstration
Method section.

\section{Method}
\subsection{Participants}
We had a lot of people in this study, show in Figure \ref{fig:Figure1}.

\begin{figure}[!htp]
    \caption{This is my figure caption.}
    \includegraphics[bb=0in 0in 2.5in 2.5in, height=1.4in, width=1.4in]{Figure1.pdf}
    \foreignlanguage{hungarian}{\figurenote{This is an awesome figure with a Hungarian note label.}}
    \label{fig:Figure1}
\end{figure}

\subsection{Materials}
Several materials were used for this project.  Some of them were
already created for prior research.

\subsubsection{Paper-and-Pencil Instrument}
We used an instrument that we found to be highly successful.

\paragraph{Reliability}
The reliability of this instrument is extraordinary.

\paragraph{Validity}
We now discuss the validity of our instrument.

\subparagraph{Face validity} The face validity is exceptionally
strong.  Everyone should be impressed.

\subparagraph{Construct validity} Also very strong.

\subsection{Design}
This section describes the study's design.

\subsection{Procedure}
The procedure was fairly straightforward, yet required
attention to detail.

\section{Results}
Table \ref{tab:ComplexTable} contains some sample data.  Our
statistical prowess in analyzing these data is unmatched.

\begin{table}[htbp]
  \vspace*{2em}
  \begin{threeparttable}
    \caption{A Complex Table}
    \label{tab:ComplexTable}
    \begin{tabular}{@{}lrrr@{}}         \toprule
    Distribution type  & \multicolumn{2}{l}{Percentage of} & Total number   \\
                       & \multicolumn{2}{l}{targets with}  & of trials per  \\
                       & \multicolumn{2}{l}{segment in}    & participant    \\ \cmidrule(r){2-3}
                                    &  Onset  &  Coda            &          \\ \midrule
    Categorical -- onset\tabfnm{a}  &    100  &     0            &  196     \\
    Probabilistic                   &     80  &    20\tabfnm{*}  &  200     \\
    Categorical -- coda\tabfnm{b}   &      0  &   100\tabfnm{*}  &  196     \\ \midrule
    \end{tabular}
    \tablenote{All data are approximate.

            \tabfnt{a}Categorical may be onset.
            \tabfnt{b}Categorical may also be coda.

            \tabfnt{*}\textit{p} < .05.
            \tabfnt{**}\textit{p} < .01.
         }
  \end{threeparttable}
\end{table}

\section{Discussion}
This is a lengthy and erudite discussion.  It demonstrates amazing
skill in interpreting the results for the masses.

\printbibliography

\end{document}

%</shortsample>
%    \end{macrocode}
%\end{macro}
%
%\begin{macro}{longsample.tex}
%    \begin{macrocode}
%<*longsample>
\documentclass[man]{apa7}

\usepackage{lipsum}

\usepackage[dutch,american]{babel}

\usepackage{csquotes}
\usepackage[style=apa,backend=biber]{biblatex}
\addbibresource{bibliography.bib}

\title{Sample APA-Style Document Using the \textsf{apa7} Package}
\shorttitle{Sample Document}

\authorsnames{Daniel A. Weiss, J. Random Hacker}
\authorsaffiliations{A University Somewhere}

\leftheader{Weiss}

\abstract{\lipsum[1]}

\keywords{APA style, demonstration}

\authornote{
   \addORCIDlink{Daniel A. Weiss}{0000-0000-0000-0000}

  Correspondence concerning this article should be addressed to Daniel A. Weiss, Department of Educational Psychology, Counseling and
  Special Education, A University Somewhere, 123 Main St., Oneonta, NY
  13820.  E-mail: daniel.weiss.led@gmail.com}


\begin{document}
\maketitle
\lipsum[2]

\Textcite{vonDavier2011} said this,
too \parencite{vonDavier2011,Lassen2006}.  Further evidence comes from
other sources \parencite{Shotton1989,Lassen2006}.  \lipsum[3]

\section{Method}
\subsection{Participants}
\lipsum[4]

\subsection{Materials}
\lipsum[5]

\subsection{Design}
\lipsum[6]

\subsection{Procedure}
\lipsum[7]

\subsubsection{Instrument \#1}
\lipsum[8]

\paragraph{Reliability}
\lipsum[9]

\subparagraph{Inter-rater reliability}
\lipsum[10]

\subparagraph{Test-retest reliability}
\lipsum[11]

\paragraph{Validity}
\lipsum[12]

\subparagraph{Face validity}
\lipsum[13]

\subparagraph{Construct validity}
\lipsum[14]


\section{Results}
Table~\ref{tab:BasicTable} summarizes the data. \lipsum[15]

\begin{table}
  \caption{Sample Basic Table}
  \label{tab:BasicTable}
  \begin{tabular}{@{}llr@{}}         \toprule
  \multicolumn{2}{c}{Item}        \\ \cmidrule(r){1-2}
  Animal    & Description & Price \\ \midrule
  Gnat      & per gram    & 13.65 \\
            & each        &  0.01 \\
  Gnu       & stuffed     & 92.50 \\
  Emu       & stuffed     & 33.33 \\
  Armadillo & frozen      &  8.99 \\ \bottomrule
  \end{tabular}
\end{table}

\begin{figure}
    \caption{This is my first figure caption.}
    \includegraphics[bb=0in 0in 2.5in 2.5in, height=2.5in, width=2.5in]{Figure1.pdf}
    \foreignlanguage{dutch}figurenote{This is a great figure with a Dutch note label.}
    \label{fig:Figure1}
\end{figure}

Figure~\ref{fig:Figure1} shows this trend. \lipsum[16]

\section{Discussion}
\lipsum[17]

\lipsum[18]

\lipsum[19]


\printbibliography

\appendix

\section{Instrument}
\label{app:instrument}

As shown in Figure~\ref{fig:Figure2}, these results are impressive. \lipsum[20]

\begin{figure}
    \caption{This is my second figure caption.}
    \includegraphics[bb=0in 0in 2.5in 2.5in, height=2.5in, width=2.5in]{Figure1.pdf}
    \figurenote{This is a great appendix figure.}
    \label{fig:Figure2}
\end{figure}

\lipsum[21]
\section{Pilot Data}
\label{app:surveydata}

The detailed results are shown in Table~\ref{tab:DeckedTable}. \lipsum[22]

\begin{table}
  \begin{threeparttable}
    \caption{A More Complex Decked Table}
    \label{tab:DeckedTable}
    \begin{tabular}{@{}lrrr@{}}         \toprule
    Distribution type  & \multicolumn{2}{l}{Percentage of} & Total number   \\
                       & \multicolumn{2}{l}{targets with}  & of trials per  \\
                       & \multicolumn{2}{l}{segment in}    & participant    \\ \cmidrule(r){2-3}
                                    &  Onset  &  Coda            &          \\ \midrule
    Categorical -- onset\tabfnm{a}  &    100  &     0            &  196     \\
    Probabilistic                   &     80  &    20\tabfnm{*}  &  200     \\
    Categorical -- coda\tabfnm{b}   &      0  &   100\tabfnm{*}  &  196     \\ \midrule
    \end{tabular}
    \tablenote{All data are approximate.

            \tabfnt{a}Categorical may be onset.
            \tabfnt{b}Categorical may also be coda.

            \tabfnt{*}\textit{p} < .05.
            \tabfnt{**}\textit{p} < .01.
         }
  \end{threeparttable}
\end{table}

\lipsum[23]

\end{document}

%</longsample>
%    \end{macrocode}
%\end{macro}
%
% %%%%%%%%%%%  CONFIGURATION FILES  %%%%%%%%%%%%
%
%\begin{macro}{apa7-american.trsl}
%    \begin{macrocode}
%<*american>
\ProvideDictionaryFor{american}{apa7}[2022/12/03 v1.00 apa7 configuration for American English]

\NewDictTranslation{lastauthorseparator}{and}
\NewDictTranslation{acksname}{Author Note}
\NewDictTranslation{keywordname}{Keywords}
\NewDictTranslation{notesname}{Footnotes}
\NewDictTranslation{notelabel}{Note}

%</american>
%    \end{macrocode}
%\end{macro}
%
%
%\begin{macro}{apa7-british.trsl}
%    \begin{macrocode}
%<*british>
\ProvideDictionaryFor{british}{apa7}[2022/12/03 v1.00 apa7 configuration for British English]

\NewDictTranslation{lastauthorseparator}{and}
\NewDictTranslation{acksname}{Author Note}
\NewDictTranslation{keywordname}{Keywords}
\NewDictTranslation{notesname}{Footnotes}
\NewDictTranslation{notelabel}{Note}

%</british>
%    \end{macrocode}
%\end{macro}
%
%
%\begin{macro}{apa7-english.trsl}
%    \begin{macrocode}
%<*english>
\ProvideDictionaryFor{english}{apa7}[2022/12/03 v1.00 apa7 configuration for non-American English]

\NewDictTranslation{lastauthorseparator}{and}
\NewDictTranslation{acksname}{Author Note}
\NewDictTranslation{keywordname}{Keywords}
\NewDictTranslation{notesname}{Footnotes}
\NewDictTranslation{notelabel}{Note}

%</english>
%    \end{macrocode}
%\end{macro}
%
%
%\begin{macro}{apa7-french.trsl}
%    \begin{macrocode}
%<*french>
\ProvideDictionaryFor{french}{apa7}[2022/12/03 v1.00 apa7 configuration for French] % translation by lejun[at]gmx.fr

\NewDictTranslation{lastauthorseparator}{et}% and
\NewDictTranslation{acksname}{Note de l'auteur}% Author Note
\NewDictTranslation{keywordname}{Mots Clés}% Keywords
\NewDictTranslation{notesname}{Notes de bas de page}% Footnotes
\NewDictTranslation{notelabel}{Note}%Note

%</french>
%    \end{macrocode}
%\end{macro}
%
%
%\begin{macro}{apa7-dutch.trsl}
%    \begin{macrocode}
%<*dutch>
\ProvideDictionaryFor{dutch}{apa7}[2022/12/03 v1.00 apa7 configuration for Dutch]

\NewDictTranslation{lastauthorseparator}{en}% and
\NewDictTranslation{acksname}{Auteursnoot}% Author Note
\NewDictTranslation{keywordname}{Trefwoorden}% Keywords
\NewDictTranslation{notesname}{Voetnoten}% Footnotes
\NewDictTranslation{notelabel}{Opmerking}%Note

%</dutch>
%    \end{macrocode}
%\end{macro}
%
%
%\begin{macro}{apa7-german.trsl}
%    \begin{macrocode}
%<*german>
\ProvideDictionaryFor{german}{apa7}[2022/12/03 v1.00 apa7 configuration for German]

\NewDictTranslation{lastauthorseparator}{und}% and
\NewDictTranslation{acksname}{Autorenhinweis}% Author Note
\NewDictTranslation{keywordname}{Schl\"usselw\"orter}% Keywords
\NewDictTranslation{notesname}{Fußnoten}% Footnotes
\NewDictTranslation{notelabel}{Anmerkung}%Note

%</german>
%    \end{macrocode}
%\end{macro}
%
%
%\begin{macro}{apa7-ngerman.trsl}
%    \begin{macrocode}
%<*ngerman>
\ProvideDictionaryFor{ngerman}{apa7}[2022/12/03 v1.00 apa7 configuration for "new" German]

\NewDictTranslation{lastauthorseparator}{und}% and
\NewDictTranslation{acksname}{Autorenhinweis}% Author Note
\NewDictTranslation{keywordname}{Schl\"usselw\"orter}% Keywords
\NewDictTranslation{notesname}{Fußnoten}% Footnotes
\NewDictTranslation{notelabel}{Anmerkung}%Note

%</ngerman>
%    \end{macrocode}
%\end{macro}
%
%
%\begin{macro}{apa7-greek.trsl}
%    \begin{macrocode}
%<*greek>
\ProvideDictionaryFor{greek}{apa7}[2022/12/03 v1.00 apa7 configuration for Greek]

\NewDictTranslation{lastauthorseparator}{και}% and
\NewDictTranslation{acksname}{Σημείωση του συγγραφέα}% Author Note
\NewDictTranslation{keywordname}{λέξεις-κλειδιά}% Keywords
\NewDictTranslation{notesname}{υποσημείωση}% Footnotes
\NewDictTranslation{notelabel}{Σημείωση}%Note

%</greek>
%    \end{macrocode}
%\end{macro}
%
%
%\begin{macro}{apa7-czech.trsl}
%    \begin{macrocode}
%<*czech>
\ProvideDictionaryFor{czech}{apa7}[2022/12/03 v1.00 apa7 configuration for Czech]

\NewDictTranslation{lastauthorseparator}{a}% and
\NewDictTranslation{acksname}{Poznámka~autora}% Author Note
\NewDictTranslation{keywordname}{Klícová~1slova}% Keywords
\NewDictTranslation{notesname}{Poznámky}% Footnotes
\NewDictTranslation{notelabel}{Nota}%Note

%</czech>
%    \end{macrocode}
%\end{macro}
%
%
%\begin{macro}{apa7-turkish.trsl}
%    \begin{macrocode}
%<*turkish>
\ProvideDictionaryFor{turkish}{apa7}[2022/12/03 v1.00 apa7 configuration for Turkish]

\NewDictTranslation{lastauthorseparator}{ve}% and
\NewDictTranslation{acksname}{Yazar Notu}% Author Note
\NewDictTranslation{keywordname}{Anahtar Kelimeler}% Keywords
\NewDictTranslation{notesname}{Dipnotlar}% Footnotes
\NewDictTranslation{notelabel}{Not}%Note

%</turkish>
%    \end{macrocode}
%\end{macro}
%
%
%\begin{macro}{apa7-spanish.trsl}
%    \begin{macrocode}
%<*spanish>
\ProvideDictionaryFor{spanish}{apa7}[2022/12/03 v1.00 apa7 configuration for Spanish]

\NewDictTranslation{lastauthorseparator}{y}% and
\NewDictTranslation{acksname}{Nota del Autor}% Author Note
\NewDictTranslation{keywordname}{Palabras Claves}% Keywords
\NewDictTranslation{notesname}{Notas}% Footnotes
%Note

%</spanish>
%    \end{macrocode}
%\end{macro}
%
%
%\begin{macro}{apa7-hungarian.trsl}
%    \begin{macrocode}
%<*hungarian>
\ProvideDictionaryFor{hungarian}{apa7}[2022/12/03 v1.00 apa7 configuration for Hungarian - added by Robert Fodor]

\NewDictTranslation{lastauthorseparator}{és}% and
\NewDictTranslation{acksname}{Szerzői megjegyzés}% Author Note
\NewDictTranslation{keywordname}{Kulcsszavak}% Keywords
\NewDictTranslation{notesname}{Lábjegyzet}% Footnotes
\NewDictTranslation{notelabel}{Megjegyzés}% Note
%Note

%</hungarian>
%    \end{macrocode}
%\end{macro}
%
%
%\begin{macro}{APA7endfloat.cfg}
%    \begin{macrocode}
%<*APAendfloat>
\ProvidesFile{APA7endfloat.cfg}[2021/03/03 v2.00 apa7 configuration endfloat]
%%
%% This is file `endfloat.cfg',
%% modifed from the original supplied efxmpl.cfg with the endfloat package
%% to handle both sideways floats and longtable
%%
%% Description: LaTeX style to put figures and tables at end of article
%% Keywords: LaTeX, style-option, float, figure, table
%% Authors: James Darrell McCauley <jdm5548@diamond.tamu.edu>,
%%          Jeff Goldberg <jeffrey@goldmark.org>,
%%          Axel Sommerfeldt <axel.sommerfeldt@f-m.fm>
%% Maintainer: -none-
%% Latest Version: Version 2.7 <April 2019>
%%
\RequirePackage{rotating}
\RequirePackage{longtable}
\DeclareDelayedFloatFlavor{sidewaysfigure}{figure}
\DeclareDelayedFloatFlavor{sidewaystable}{table}
\DeclareDelayedFloatFlavor{longtable}{table}
\endinput%</APAendfloat>
%    \end{macrocode}
%\end{macro}
%
%
%\Finale
